\documentclass[11pt]{article}
\usepackage[margin=1.1in]{geometry}
\usepackage{amsmath,amssymb,amsthm,mathtools,booktabs,array,tabularx,longtable}
\usepackage[dvipsnames]{xcolor}
\usepackage[colorlinks=true,linkcolor=MidnightBlue,citecolor=MidnightBlue,urlcolor=MidnightBlue]{hyperref}
\usepackage{fancyhdr}
\pagestyle{fancy}\fancyhf{}\renewcommand{\headrulewidth}{0pt}
\fancyfoot[R]{\thepage}
\fancyfoot[L]{\scriptsize\color{gray}Generated by the Transport Lab from approved parts of its store; not a deposited paper. Papers of record: \href{https://zenodo.org/communities/transport-papers}{Zenodo, Transport Papers}}
\fancypagestyle{plain}{\fancyhf{}\fancyfoot[R]{\thepage}\fancyfoot[L]{\scriptsize\color{gray}Generated by the Transport Lab from approved parts of its store; not a deposited paper. Papers of record: \href{https://zenodo.org/communities/transport-papers}{Zenodo, Transport Papers}}\renewcommand{\headrulewidth}{0pt}}
\newtheorem{theorem}{Theorem}[section]
\newtheorem{proposition}[theorem]{Proposition}
\newtheorem{lemma}[theorem]{Lemma}
\newtheorem{corollary}[theorem]{Corollary}
\newtheorem{conjecture}[theorem]{Conjecture}
\newtheorem{refutation}[theorem]{Refutation}
\newtheorem{prediction}[theorem]{Prediction}
\newtheorem{principle}[theorem]{Principle}
\newtheorem{observation}[theorem]{Observation}
\newtheorem{axiom}[theorem]{Axiom}
\theoremstyle{definition}
\newtheorem{definition}[theorem]{Definition}
\newtheorem{fact}[theorem]{Fact}
\newtheorem{assembled}[theorem]{Assembly}
\newtheorem{postulate}[theorem]{Postulate}
\newtheorem{premise}[theorem]{Premise}
\newtheorem{hypothesis}[theorem]{Hypothesis}
\newtheorem{remark}[theorem]{Remark}
\newtheorem{example}[theorem]{Example}
\theoremstyle{plain}
\setlength{\marginparwidth}{0.8in}
\newcommand{\changedmark}[1]{\marginpar{\raggedright\scriptsize\color{BrickRed}#1}}
\providecommand{\Aa}{}\renewcommand{\Aa}{\mathcal{A}}
\providecommand{\Ad}{}\renewcommand{\Ad}{\widetilde{\mathrm{Ad}}}
\providecommand{\Aev}{}\renewcommand{\Aev}{\mathcal{A}^{+}}
\providecommand{\Aod}{}\renewcommand{\Aod}{\mathcal{A}^{-}}
\providecommand{\B}{}\renewcommand{\B}{\mathrm{B}}
\providecommand{\C}{}\renewcommand{\C}{\mathbb{C}}
\providecommand{\CAR}{}\renewcommand{\CAR}{\mathrm{CAR}}
\providecommand{\CS}{}\renewcommand{\CS}{{\mathbb{S}^{+}}}
\providecommand{\CSm}{}\renewcommand{\CSm}{{\mathbb{S}^{-}}}
\providecommand{\Cl}{}\renewcommand{\Cl}{\mathrm{Cl}}
\providecommand{\Con}{}\renewcommand{\Con}{\mathcal{C}}
\providecommand{\Dl}{}\renewcommand{\Dl}{\Delta}
\providecommand{\EN}{}\renewcommand{\EN}{E_{N}}
\providecommand{\End}{}\renewcommand{\End}{\mathrm{End}}
\providecommand{\F}{}\renewcommand{\F}{\mathbb{F}}
\providecommand{\G}{}\renewcommand{\G}{\Gamma}
\providecommand{\HH}{}\renewcommand{\HH}{\mathbb{H}}
\providecommand{\IS}{}\renewcommand{\IS}{I\sigma_3}
\providecommand{\J}{}\renewcommand{\J}{\mathbf{J}}
\providecommand{\LH}{}\renewcommand{\LH}{L_{H}}
\providecommand{\Mat}{}\renewcommand{\Mat}{\mathrm{M}}
\providecommand{\Mo}{}\renewcommand{\Mo}{M^{\circ}}
\providecommand{\Mp}{}\renewcommand{\Mp}{M^{+}}
\providecommand{\Msun}{}\renewcommand{\Msun}{M_{\odot}}
\providecommand{\Mt}{}\renewcommand{\Mt}{\widetilde{M}}
\providecommand{\Ncone}{}\renewcommand{\Ncone}{\mathcal{N}}
\providecommand{\Neff}{}\renewcommand{\Neff}{N_{\mathrm{eff}}}
\providecommand{\PCS}{}\renewcommand{\PCS}{\mathbb{P}{\mathbb{S}^{+}}}
\providecommand{\PT}{}\renewcommand{\PT}{\mathbb{PT}}
\providecommand{\Pin}{}\renewcommand{\Pin}{\mathrm{Pin}}
\providecommand{\Q}{}\renewcommand{\Q}{\mathbb{Q}}
\providecommand{\R}{}\renewcommand{\R}{\mathbb{R}}
\providecommand{\Ric}{}\renewcommand{\Ric}{\mathrm{Ric}}
\providecommand{\SBH}{}\renewcommand{\SBH}{S_{\mathrm{BH}}}
\providecommand{\Sc}{}\renewcommand{\Sc}{\mathbb{S}}
\providecommand{\Sch}{}\renewcommand{\Sch}{\mathrm{P}}
\providecommand{\Sig}{}\renewcommand{\Sig}{\Sigma}
\providecommand{\Som}{}\renewcommand{\Som}{S_\omega}
\providecommand{\Sp}{}\renewcommand{\Sp}{\mathbb{S}}
\providecommand{\Sph}{}\renewcommand{\Sph}{\mathrm{S}}
\providecommand{\Spi}{}\renewcommand{\Spi}{S_\pi}
\providecommand{\Spin}{}\renewcommand{\Spin}{\mathrm{Spin}}
\providecommand{\Sth}{}\renewcommand{\Sth}{S_{\mathrm{th}}}
\providecommand{\T}{}\renewcommand{\T}{\mathbb{T}}
\providecommand{\TH}{}\renewcommand{\TH}{T_{\mathrm{H}}}
\providecommand{\Tr}{}\renewcommand{\Tr}{\operatorname{Tr}}
\providecommand{\UHF}{}\renewcommand{\UHF}{\mathrm{UHF}}
\providecommand{\Z}{}\renewcommand{\Z}{\mathbb{Z}}
\providecommand{\Zc}{}\renewcommand{\Zc}{\mathrm{Z}}
\providecommand{\cmg}{}\renewcommand{\cmg}{\mathrm{cm^{2}\,g^{-1}}}
\providecommand{\dS}{}\renewcommand{\dS}{\mathrm{dS}}
\providecommand{\dsc}{}\renewcommand{\dsc}{\mathrm{dsc}}
\providecommand{\dtPG}{}\renewcommand{\dtPG}{dt}
\providecommand{\dtS}{}\renewcommand{\dtS}{dt_{\mathrm{S}}}
\providecommand{\ellv}{}\renewcommand{\ellv}{\ell}
\providecommand{\er}{}\renewcommand{\er}{e_r}
\providecommand{\et}{}\renewcommand{\et}{e_t}
\providecommand{\ev}{}\renewcommand{\ev}{\mathrm{ev}}
\providecommand{\gl}{}\renewcommand{\gl}{\mathfrak{gl}}
\providecommand{\hatV}{}\renewcommand{\hatV}{\widehat{V}}
\providecommand{\hb}{}\renewcommand{\hb}{\bar{h}}
\providecommand{\hf}{}\renewcommand{\hf}{\hat{f}}
\providecommand{\id}{}\renewcommand{\id}{\mathrm{id}}
\providecommand{\kpc}{}\renewcommand{\kpc}{\mathrm{kpc}}
\providecommand{\lP}{}\renewcommand{\lP}{\ell_{P}}
\providecommand{\mD}{}\renewcommand{\mD}{\mathbf{m}_{D}}
\providecommand{\mb}{}\renewcommand{\mb}{\bar m}
\providecommand{\mnu}{}\renewcommand{\mnu}{\Sigma m_{\nu}}
\providecommand{\nInf}{}\renewcommand{\nInf}{n_\infty}
\providecommand{\nO}{}\renewcommand{\nO}{n_0}
\providecommand{\ninf}{}\renewcommand{\ninf}{n_{\infty}}
\providecommand{\nz}{}\renewcommand{\nz}{n_{0}}
\providecommand{\proj}{}\renewcommand{\proj}{\operatorname{proj}}
\providecommand{\pv}{}\renewcommand{\pv}{\pi}
\providecommand{\rad}{}\renewcommand{\rad}{\mathrm{rad}}
\providecommand{\sgn}{}\renewcommand{\sgn}{\operatorname{sgn}}
\providecommand{\shA}{}\renewcommand{\shA}{\mathcal{A}}
\providecommand{\shC}{}\renewcommand{\shC}{\mathcal{C}}
\providecommand{\shE}{}\renewcommand{\shE}{\mathcal{E}}
\providecommand{\shJ}{}\renewcommand{\shJ}{\mathcal{J}}
\providecommand{\shK}{}\renewcommand{\shK}{\mathcal{K}}
\providecommand{\shL}{}\renewcommand{\shL}{L}
\providecommand{\shW}{}\renewcommand{\shW}{\mathcal{W}}
\providecommand{\sll}{}\renewcommand{\sll}{\mathfrak{sl}}
\providecommand{\so}{}\renewcommand{\so}{\mathfrak{so}}
\providecommand{\su}{}\renewcommand{\su}{\mathfrak{su}}
\providecommand{\supp}{}\renewcommand{\supp}{\operatorname{supp}}
\providecommand{\tr}{}\renewcommand{\tr}{\operatorname{tr}}
\providecommand{\uu}{}\renewcommand{\uu}{\mathfrak{u}}
\providecommand{\br}{}\renewcommand{\br}[1]{\langle #1\rangle}
\providecommand{\bra}{}\renewcommand{\bra}[1]{\langle #1|}
\providecommand{\file}{}\renewcommand{\file}[1]{\texttt{#1}}
\providecommand{\hp}{}\renewcommand{\hp}[2]{( #1 \,|\, #2 )}
\providecommand{\ip}{}\renewcommand{\ip}[2]{\langle #1, #2\rangle}
\providecommand{\ket}{}\renewcommand{\ket}[1]{|#1\rangle}
\providecommand{\rel}{}\renewcommand{\rel}[2]{D\bigl(#1 \,\|\, #2\bigr)}
\providecommand{\rev}{}\renewcommand{\rev}[1]{\widetilde{#1}}
\providecommand{\slashed}{}\renewcommand{\slashed}[1]{{#1}\mkern-10mu/}
\title{\bfseries Where the Pairing Generator Is: The connection moves the vacuum by the optical scalars of its null direction: Companion note to T5 and T16: a dictionary, one claim closed for the pre-ascent fibre, one not closed}
\author{Leon Butler\\ \normalsize\href{mailto:lbutler.math@gmail.com}{\texttt{lbutler.math@gmail.com}}\\ \normalsize ORCID: \href{https://orcid.org/0009-0006-4449-5617}{0009-0006-4449-5617}}
\date{Generated by the Transport Lab, 27 October 2026}
\begin{document}
\maketitle
\begin{abstract}
\normalsize
The singularity paper places the pairing terms, the generators that change a Witt decomposition and so create entanglement from a vacuum, in the tangent of the vacuum orbit, and records two things it does not do: it compares two Witt decompositions without deriving them, and supplies no generator in that tangent. This note states both as claims with kill conditions fixed before computation, on the laboratory paper's pre-ascent Dirac fibre, two modes at a point, not the boundary-mode line. For the vacuum given by a null direction $\ell$ and a chirality bit, the stabiliser in $\so(1,3)$ is the parabolic algebra of $\ell$, and the functional $\pv(B) = \ip{B\cdot\ell}{m}$ that cuts it out is, on the Fock module, the pair-creation amplitude $-\bra{12}B\ket{0}$. The connection therefore has a pairing part, the transverse derivative of the null direction: the Newman-Penrose coefficients $\kappa, \tau, \sigma, \rho$. What is proved is the dictionary and the exact Schwarzschild classes; the identifications are imports, and the whole stands without the transport principle. In Schwarzschild, with $\ell$ the outgoing radial direction, only $\rho$ survives: zero along the $t$- and $r$-lines, zero on the horizon in a regular frame, opposite in sign on the two sides. Holonomy moves the vacuum only for loops with an $n \wedge \mb$ component, $\ell$ being a repeated principal null direction. The mismatch angle is not reached: the two frames share a vacuum, the selected transport has zero class, and no holonomy carries a mode weight. Nothing here pairs the ladders at two points.
\end{abstract}
\medskip
\noindent\textbf{MSC 2020.} Primary 15A66; Secondary 83C57, 81R25, 53C27.\\
\textbf{Keywords.} Clifford algebras; Witt decomposition; pure spinors; Newman-Penrose spin coefficients; Schwarzschild; Bogoliubov transformation; transport principle.
\newpage
\noindent \textbf{Authorship and disclosure. }The founding idea of the series is the author's: that division by zero in a Clifford algebra should be read not as a prohibition but as a transport into a higher-dimensional algebra, and that the transported division need not be accompanied by an inverse operation. The direction of investigation and the critical review of the results are likewise the author's. The mathematical development and the drafting of the text were carried out in extended collaboration with Claude Fable 5 (Anthropic), an AI system.

\noindent\fbox{\parbox{\dimexpr\linewidth-2\fboxsep-2\fboxrule}{\textbf{This is not a deposited paper.} It was generated by the Transport Lab on 27 October 2026 and assembled from approved parts of its store. It follows results across several papers of the series; each is stated with its proof and credited to the paper that first proved it. The official versions of those papers are on Zenodo, in the Transport Papers community, \url{https://zenodo.org/communities/transport-papers}.}}
\newpage

\section{Setting, and what is imported}

\phantomsection\label{c:note.pairing.setting-setting-what-imported}
The singularity paper reads entanglement as a mismatch of Witt decompositions and locates the generators that can produce it \cite[Prop.~13.4]{ButlerT5}: on the Fock space of $\nu$ modes the bivectors of $\so(2\nu)$ split into hops and pair annihilations, which stabilise the vacuum ray, and pair creations, which span the tangent of its orbit. The same paper then records, in one remark, the two things it does not do \cite[Rem.~13.5]{ButlerT5}: the framework compares two Witt decompositions and does not derive them, and the derivation of the one pairing angle its chain uses, the angle of its regularity premise, lies outside its reach. Its locality premise (L$_E$) \cite[Prem.~7.30]{ButlerT5} excludes pairing terms on the ground that the completion supplies none.

This note asks where, in the framework's own data, a pairing generator would have to come from, and answers: in the connection. The laboratory paper says in words that the connection carries the vacuum's bivector from point to point and returns it tilted \cite[Prop.~4.1(3) and the paragraph following it]{ButlerT16}. Stated exactly, the tilt has a stabiliser part and a pairing part, and the pairing part is the set of Newman-Penrose spin coefficients that measure how the null direction fails to be parallel. The note computes it, in the series' own conventions and in exact arithmetic, and then applies the kill conditions of the section ``The two claims, and the kill conditions fixed before computation'' to the two claims.

Three things are said on this page because the rest depends on them. First, the space. The singularity paper's premise governs the Fock space of its boundary modes, one ladder pair attached to each point of the exceptional set $E$ \cite[Prem.~7.30, Fact~5.9]{ButlerT5}, on which a pairing $v^{*}_{x}v^{*}_{y}$ is cross-point by construction and a single point carries no pairing at all. The fibre of this note is a different space: the pre-ascent Dirac spinor at an exterior point, read as two fermionic modes, the complex capacity two of the laboratory paper \cite[Prop.~3.2(1)]{ButlerT16}. What is found here is found on that fibre, and the section ``What this does to the two claims'' says what it does and does not imply for the premise. Second, the risk. The positive claim's kill conditions could not have fired (the section ``What this does to the two claims''); the registration fixed the computation and its conventions in advance, and the note's content is a dictionary between the series' objects and known ones, not a derivation that might have failed. Third, the vocabulary, for a reader of general relativity and geometric algebra who has not read the series: the \emph{ascent} is the adjunction of a hyperbolic plane $(1,1)$ that the transport principle forces where division fails \cite[Thm.~4.13]{ButlerT1}, the \emph{completion} is its result, the \emph{exceptional set} $E$ is the locus where it happens \cite[Thm.~8.12]{ButlerT4}, the \emph{descent} is the one-way map back \cite[Thm.~7.3]{ButlerT5}, the \emph{flag bit} is the orientation choosing which half of a null pair creates \cite[Def.~2.10]{ButlerT14}, and the two-ledger words \emph{transport-essential}, \emph{convenient} and \emph{absent} grade a result by whether the $(1,1)$ increment does irreplaceable work in it, some work, or none \cite{ButlerT17}.

Notation follows the series. $\Cl(1,3)$ has $e_0^2 = 1$ and $e_i^2 = -1$; $\ip{a}{b} = \tfrac12(ab + ba)$ is the scalar product of vectors, extended bilinearly (not sesquilinearly) to $\C^4$; a bivector $B$ acts on vectors by $B\cdot v = \tfrac12[B, v]$, which is the $\so(1,3)$ action. The null tetrad is
\[
\ell = e_0 + e_1, \qquad n = e_0 - e_1, \qquad m = e_2 + i e_3, \qquad \mb = e_2 - i e_3,
\]
with $\ip{\ell}{n} = 2$, $\ip{m}{\mb} = -2$ and all other pairings zero. The framework's vacuum at a point is the maximal isotropic subspace $W = \operatorname{span}\{\ell, m\}$ of $\C^4$: $\ell$ is the radial null direction of the descent, the ladder vector of the singularity paper's Step 3 up to the flag bit \cite[Sec.~4, Lem.~4.1]{ButlerT5}, and the sign of $i$ in $m$ is the chirality bit. The Fock module is that of two modes with annihilators $a_1 = c(\ell)/2$, $a_2 = c(m)/2$ and creators $a_1^{\dagger} = c(n)/2$, $a_2^{\dagger} = -c(\mb)/2$, so that $e_0 = a_1 + a_1^{\dagger}$, $e_1 = a_1 - a_1^{\dagger}$, $e_2 = a_2 - a_2^{\dagger}$, $e_3 = -i(a_2 + a_2^{\dagger})$; this is the two-mode reading of the Dirac spinor, the complex capacity two of the laboratory paper.

\begin{fact}[Imported: the stabiliser and the tangent (external ingredients)]\label{c:pairing.imported-stabiliser-tangent-external}
The results from outside the series on which \ref{c:pairing.imported-stabiliser-tangent} draws, as cited there: 'E. Cartan, Lecons sur la th'eorie des spineurs, Hermann, Paris (1938).; C. Chevalley, The Algebraic Theory of Spinors, Columbia University Press (1954).; D. J. Thouless, Stability conditions and nuclear rotations in the Hartree-Fock theory, Nucl. Phys. 21 (1960) 225--232..
\emph{Source:} 'E. Cartan, Lecons sur la th'eorie des spineurs, Hermann, Paris (1938).; C. Chevalley, The Algebraic Theory of Spinors, Columbia University Press (1954).; D. J. Thouless, Stability conditions and nuclear rotations in the Hartree-Fock theory, Nucl. Phys. 21 (1960) 225--232..
\end{fact}

\begin{assembled}[Imported: the stabiliser and the tangent]\label{c:pairing.imported-stabiliser-tangent}
On the Fock space of $\nu \geq 2$ boundary modes the bivectors of $\so(2\nu)$ split as hops $a_i^{\dagger}a_j$ (with the Cartan elements $a_i^{\dagger}a_i - \tfrac12$), pair annihilations $a_ia_j$ and pair creations $a_i^{\dagger}a_j^{\dagger}$. On the vacuum every hop with $i \neq j$ and every annihilation vanishes and every Cartan element is a scalar, so hops and annihilations are the stabiliser of the vacuum ray; the creations span the tangent of the orbit of Witt vacua, of dimension $\nu(\nu-1)/2$. A hop rotation fixes the vacuum, a pairing rotation moves it, and the necessity witness of the singularity paper's envelope corollary, $\ket{00} \mapsto \tfrac35\ket{00} + \tfrac45\ket{11}$, is the pairing rotation on one pair at $(\cos\theta, \sin\theta) = (\tfrac35, \tfrac45)$. \emph{Provenance:} \cite[Prop.~13.4, Cor.~11.6]{ButlerT5}; the orbit is the spinor variety of \cite{Cartan1938, Chevalley1954}, imported there, and the exponential form of a rotated vacuum is Thouless's \cite{Thouless1960}, imported by the laboratory paper \cite[Fact~2.10]{ButlerT16}.
\end{assembled}
\begin{proof}
On the Fock space of $\nu \geq 2$ boundary modes the bivectors of $\so(2\nu)$ split as hops $a_i^{\dagger}a_j$ (with the Cartan elements $a_i^{\dagger}a_i - \tfrac12$), pair annihilations $a_ia_j$ and pair creations $a_i^{\dagger}a_j^{\dagger}$. On the vacuum every hop with $i \neq j$ and every annihilation vanishes and every Cartan element is a scalar, so hops and annihilations are the stabiliser of the vacuum ray; the creations span the tangent of the orbit of Witt vacua, of dimension $\nu(\nu-1)/2$. A hop rotation fixes the vacuum, a pairing rotation moves it, and the necessity witness of the singularity paper's envelope corollary, $\ket{00} \mapsto \tfrac35\ket{00} + \tfrac45\ket{11}$, is the pairing rotation on one pair at $(\cos\theta, \sin\theta) = (\tfrac35, \tfrac45)$. \emph{Provenance:} \cite[Prop.~13.4, Cor.~11.6]{ButlerT5} (\ref{c:singularity.envelope-endpoint-why-boundary-mode-locality}, \ref{c:singularity.l-stabiliser-tangent}); the orbit is the spinor variety of \cite{Cartan1938, Chevalley1954}, imported there, and the exponential form of a rotated vacuum is Thouless's \cite{Thouless1960}, imported by the laboratory paper \cite[Fact~2.10]{ButlerT16} (\ref{c:one-relation.imported-cartans-purity-theorem-thouless-form}).

This restates \ref{c:pairing.imported-stabiliser-tangent-external}.

\end{proof}

\begin{assembled}[Imported: the two readings, the capacities, and the vacuum as a null direction]\label{c:pairing.imported-two-readings-capacities-vacuum-null}
The relation $\{e_a, e_b\} = 2g_{ab}$ on a Witt basis is the canonical anticommutation relations, the Fock module being supplied by the pure-spinor construction. A real quadratic space of signature $(p,q)$ has real capacity $\min(p,q)$ and complex capacity $\lfloor (p+q)/2 \rfloor$; for $(1,3)$ these are one and two, the complex reading being the Dirac spinor as a two-mode Fock space. A vacuum at a point of spacetime is an extra datum, a null direction at the point, which the reading must supply. Every Weyl spinor of $\mathrm{Spin}(2\nu, \C)$ is pure for $2\nu \leq 6$, so at $\nu = 2$ and $\nu = 3$ every parity-definite state is the vacuum of some Witt decomposition. \emph{Provenance:} \cite[Props.~3.1, 3.2, 4.1(3), Fact~2.10]{ButlerT16}; the capacity count is \cite[Thm.~4.3]{ButlerT3}.
\end{assembled}
\begin{proof}
The relation $\{e_a, e_b\} = 2g_{ab}$ on a Witt basis is the canonical anticommutation relations, the Fock module being supplied by the pure-spinor construction. A real quadratic space of signature $(p,q)$ has real capacity $\min(p,q)$ and complex capacity $\lfloor (p+q)/2 \rfloor$; for $(1,3)$ these are one and two, the complex reading being the Dirac spinor as a two-mode Fock space. A vacuum at a point of spacetime is an extra datum, a null direction at the point, which the reading must supply. Every Weyl spinor of $\mathrm{Spin}(2\nu, \C)$ is pure for $2\nu \leq 6$, so at $\nu = 2$ and $\nu = 3$ every parity-definite state is the vacuum of some Witt decomposition. \emph{Provenance:} \cite[Props.~3.1, 3.2, 4.1(3), Fact~2.10]{ButlerT16} (\ref{c:one-relation.imported-cartans-purity-theorem-thouless-form}, \ref{c:one-relation.two-presentations}, \ref{c:one-relation.spacetime-manifold-vacuum-lines}, \ref{c:one-relation.two-capacities-three-consequences}); the capacity count is \cite[Thm.~4.3]{ButlerT3} (\ref{c:tower.witt-index-quantum-capacity}).

\end{proof}

\begin{fact}[Imported: null flags]\label{c:pairing.imported-null-flags}
A nonzero Weyl spinor $o$ of $\mathrm{Spin}(1,3)$ determines a future null direction $\ell = o\bar o$ and a flag plane through it, and the projective Weyl spinor $[o]$ is the null direction: the celestial sphere is $\mathbb{P}^1$. A rotation of the frame by an angle $\phi$ about an axis acts on the spinor by the half angle $\phi/2$. \emph{Provenance:} \cite[Sec.~1.4]{PenroseRindler1}.
\emph{Source:} R. Penrose, W. Rindler, Spinors and Space-Time, Vol. 1: Two-Spinor Calculus and Relativistic Fields, Cambridge University Press (1984)., Sec.~1.4.
\end{fact}

\begin{fact}[Imported: the spin coefficients of a null tetrad]\label{c:pairing.imported-spin-coefficients-null-tetrad}
Let $(\hb, \Omega)$ be a gravitational field on flat spacetime with covariant derivative $D$ (Definition \emph{Gauge Theory Gravity: the gravitational field on flat spacetime}), and $(\ell, n, m, \mb)$ a null tetrad of physical vector fields on its invertible region (complexified for $m$, $\mb$), with $\ip{\cdot}{\cdot}$ the flat product of physical vectors, the derived metric. The four spin coefficients of the first kind are the components on $\ell, n, m, \mb$ of the linear form $X \mapsto \ip{(X \cdot D)\, \ell}{m}$, up to the normalisation of the tetrad and the sign conventions of the formalism:
\[
\kappa \leftrightarrow \ip{(\ell \cdot D)\, \ell}{m}, \qquad \tau \leftrightarrow \ip{(n \cdot D)\, \ell}{m}, \qquad \sigma \leftrightarrow \ip{(m \cdot D)\, \ell}{m}, \qquad \rho \leftrightarrow \ip{(\mb \cdot D)\, \ell}{m}
\]
(in metric language, $\ip{\nabla_X \ell}{m}$ for the Levi-Civita derivative of the derived metric).
Geometrically: $\kappa = 0$ if and only if the integral curves of $\ell$ are geodesics (worldlines along which $\ell$ is $D$-parallel up to rescaling) (this $\kappa$ is the formalism's geodesy coefficient, not the surface gravity the singularity paper writes with the same letter); $\sigma$ is the shear of the congruence of $\ell$; $\rho$ is minus its expansion (real part) and twist (imaginary part) in the standard sign convention, so that an expanding congruence has $\rho < 0$ and the radial rays of a vanishing field ($\hb = 1$, $\Omega = 0$) have $\rho = -1/r$; and $\tau$ is the rate at which $\ell$ turns transversely along $n$. The remaining eight spin coefficients are the components of $\ip{(X \cdot D)\, \ell}{n}$, $\ip{(X \cdot D)\, m}{\mb}$ and $\ip{(X \cdot D)\, n}{\mb}$: they rescale $\ell$, rotate $m$ in its plane, and tilt $n$ and $\mb$ by null rotations about $\ell$, all of which keep $\ell$ and $m$ in their span. On a Killing horizon, the locus where a Killing field $K$ is null and normal to it, the generators have $\kappa = \sigma = \rho = 0$ in any tetrad regular there. \emph{Provenance:} \cite{NewmanPenrose1962}; \cite[Ch.~4]{PenroseRindler1}; the horizon statement \cite[Sec.~12.5]{Wald1984}.
\emph{Source:} E. Newman, R. Penrose, An approach to gravitational radiation by a method of spin coefficients, J. Math. Phys. 3 (1962) 566--578.; R. Penrose, W. Rindler, Spinors and Space-Time, Vol. 1: Two-Spinor Calculus and Relativistic Fields, Cambridge University Press (1984)., Ch.~4; R. M. Wald, General Relativity, University of Chicago Press (1984)., Sec.~12.5.
\end{fact}

\begin{assembled}[Imported: the Schwarzschild connection and curvature]\label{c:pairing.imported-schwarzschild-connection-curvature}
For the Schwarzschild gravitational field $(\hb, \Omega)$ on flat spacetime in the static gauge (Definition \emph{Gauge Theory Gravity: the gravitational field on flat spacetime}), with static orthonormal coframe $\theta^{0} = \sqrt f\,dt$, $\theta^{1} = dr/\sqrt f$, $\theta^{2} = r\,d\vartheta$, $\theta^{3} = r\sin\vartheta\,d\varphi$, $f = 1 - 2M/r$, the rotation-gauge field $\Omega$, read as the connection forms of that coframe, is $\omega^{0}{}_{1} = (f'/2\sqrt f)\,\theta^{0}$, $\omega^{2}{}_{1} = (\sqrt f/r)\,\theta^{2}$, $\omega^{3}{}_{1} = (\sqrt f/r)\,\theta^{3}$, $\omega^{3}{}_{2} = (\cot\vartheta/r)\,\theta^{3}$, and the curvature operator $\mathcal{R}(B)$ on physical bivectors (the field strength in the store's sign) is diagonal on the coordinate planes with eigenvalues $2M/r^{3}$ on the radial plane $\hat0\hat1$ and the angular plane $\hat2\hat3$ and $-M/r^{3}$ on the four mixed planes. The extreme eigenvalue distinguishes the pair $\{\sigma_{r}, I\sigma_{r}\}$ invariantly, the radial plane and its Hodge dual; the radial member is selected within the pair by the sign of the square, $\sigma_{r}^{2} = +1$ against $(I\sigma_{r})^{2} = -1$. Which of the two null directions of the radial plane is the vacuum's is the flag bit, which the series takes as initial data. In the Painlev\'e-Gullstrand gauge $\hb = 1 - \beta P$, $P(a) = (a \cdot e_{t})\,e_{r}$, the derived metric is $ds^{2} = dT^{2} - (dr + \beta\,dT)^{2} - r^{2}d\Omega^{2}$ with $\beta^{2} = 2M/r$, and $\hb$ is invertible through the horizon. \emph{Provenance:} the connection and curvature forms \cite[App.~A.2]{ButlerT5}; the pair, the sign of the square and the matching \cite[Thm.~6.1, Prop.~7.21(3)]{ButlerT5}; the Painlev\'e-Gullstrand form \cite[App.~A.1]{ButlerT5}; the flag bit \cite[Def.~2.10, Prop.~5.1]{ButlerT14}.
\end{assembled}
\begin{proof}
Read in Gauge Theory Gravity, the coframe, connection forms and curvature below are those of the Schwarzschild field in the static gauge, its rotation-gauge field and its field strength $\mathcal{R}(B)$ (Fact \ref{c:synthesis.imported-gtg-dictionary}, clauses (i), (ii) and (iv)). In the static orthonormal coframe $\theta^{0} = \sqrt f\,dt$, $\theta^{1} = dr/\sqrt f$, $\theta^{2} = r\,d\vartheta$, $\theta^{3} = r\sin\vartheta\,d\varphi$, $f = 1 - 2M/r$, the connection forms are $\omega^{0}{}_{1} = (f'/2\sqrt f)\,\theta^{0}$, $\omega^{2}{}_{1} = (\sqrt f/r)\,\theta^{2}$, $\omega^{3}{}_{1} = (\sqrt f/r)\,\theta^{3}$, $\omega^{3}{}_{2} = (\cot\vartheta/r)\,\theta^{3}$, and the curvature operator is diagonal on the coordinate planes with eigenvalues $2M/r^{3}$ on the radial plane $\hat0\hat1$ and the angular plane $\hat2\hat3$ and $-M/r^{3}$ on the four mixed planes. The extreme eigenvalue distinguishes the pair $\{\sigma_{r}, I\sigma_{r}\}$ invariantly, the radial plane and its Hodge dual; the radial member is selected within the pair by the sign of the square, $\sigma_{r}^{2} = +1$ against $(I\sigma_{r})^{2} = -1$. Which of the two null directions of the radial plane is the vacuum's is the flag bit, which the series takes as initial data. In Painlev\'e-Gullstrand form $ds^{2} = dT^{2} - (dr + \beta\,dT)^{2} - r^{2}d\Omega^{2}$ with $\beta^{2} = 2M/r$, regular through the horizon. \emph{Provenance:} the connection and curvature forms \cite[App.~A.2]{ButlerT5}; the pair, the sign of the square and the matching \cite[Thm.~6.1, Prop.~7.21(3)]{ButlerT5} (\ref{c:singularity.curvature-gauge-field-exact-shape-decoupling}, \ref{c:singularity.euclidean-rotor-where-it-lies}); the Painlev\'e-Gullstrand form \cite[App.~A.1]{ButlerT5}; the flag bit \cite[Def.~2.10, Prop.~5.1]{ButlerT14}.

\end{proof}

\begin{fact}[Imported: the flat conformal model and the incidence relation (external ingredients)]\label{c:pairing.imported-flat-conformal-model-twistor-incidence-external}
The results from outside the series on which \ref{c:pairing.imported-flat-conformal-model-twistor-incidence} draws, as cited there: T. N. Bailey, M. G. Eastwood, A. R. Gover, Thomas's structure bundle for conformal, projective and related structures, Rocky Mountain J. Math. 24 (1994) 1191--1217..
\emph{Source:} T. N. Bailey, M. G. Eastwood, A. R. Gover, Thomas's structure bundle for conformal, projective and related structures, Rocky Mountain J. Math. 24 (1994) 1191--1217..
\end{fact}

\begin{assembled}[Imported: the flat conformal model and the incidence relation]\label{c:pairing.imported-flat-conformal-model-twistor-incidence}
In $\Cl(2,4)$ with $\nz = \tfrac12(e_{+} + e_{-})$, $\ninf = e_{+} - e_{-}$, $\ip{\nz}{\ninf} = 1$, the point embedding $X(x) = \nz + x - \tfrac12 Q(x)\ninf$ is null, and the vacua at $x$ are the conformal spinors (twistors) whose isotropic three-space contains $X(x)$, a projective line of $\PCS$; two points are null-separated exactly when their lines meet. On the flat model the tractor connection is the ambient derivative, so its parallel transport is trivial, and the motion of the point's null vector $X(x)$ under a displacement is generated by the translations $\{e_a \wedge \ninf\}$, the four-dimensional nilradical of the ten-dimensional stabiliser $\so(1,3) \ltimes \R^{4}$ of $\ninf$ in $\so(2,4)$, the Poincar\'e algebra. \emph{Provenance:} \cite[Fact~2.3, Prop.~4.1(2)]{ButlerT16}; the incidence theorem and the stabiliser \cite[Thm.~4.2, Fact~4.3]{ButlerT9}; the tractor connection \cite{BEG1994}.
\end{assembled}
\begin{proof}
In $\Cl(2,4)$ with $\nz = \tfrac12(e_{+} + e_{-})$, $\ninf = e_{+} - e_{-}$, $\ip{\nz}{\ninf} = 1$, the point embedding $X(x) = \nz + x - \tfrac12 Q(x)\ninf$ is null, and the vacua at $x$ are the conformal spinors (twistors) whose isotropic three-space contains $X(x)$, a projective line of $\PCS$; two points are null-separated exactly when their lines meet. On the flat model the tractor connection is the ambient derivative, so its parallel transport is trivial, and the motion of the point's null vector $X(x)$ under a displacement is generated by the translations $\{e_a \wedge \ninf\}$, the four-dimensional nilradical of the ten-dimensional stabiliser $\so(1,3) \ltimes \R^{4}$ of $\ninf$ in $\so(2,4)$, the Poincar\'e algebra. \emph{Provenance:} \cite[Fact~2.3, Prop.~4.1(2)]{ButlerT16} (\ref{c:one-relation.imported-pure-spinors-incidence-chart}, \ref{c:one-relation.spacetime-manifold-vacuum-lines}); the incidence theorem and the stabiliser \cite[Thm.~4.2, Fact~4.3]{ButlerT9} (\ref{c:correspondence.imported-stabilisers-flat-model}, \ref{c:correspondence.incidence-relation-chart}); the tractor connection \cite{BEG1994}.

This restates \ref{c:pairing.imported-flat-conformal-model-twistor-incidence-external}.

\end{proof}

\section{The two claims, and the kill conditions fixed before computation}

\phantomsection\label{c:note.pairing.setting-two-claims-kill-conditions-fixed}
The series' method is to state the conditions under which a claim dies before computing anything. The kill conditions and predictions below are as registered, in the deposited pre-registration \cite{ButlerPreregT05pairing}; every departure between that wording and what the note does is collected in Remark \ref{c:note.pairing.departures-registered-wording}, and the assessment of how much risk each condition carried is in the section ``What this does to the two claims''. $\Q(F)$ below is the field of rational functions of $F = \sqrt f$, in which the static-frame computations are exact.

\phantomsection\label{c:note.pairing.setting-two-claims-kill-conditions-fixed-2}
\textbf{Claim R2 (the pairing generator).} The framework's own connection, relative to the framework's own vacuum field, has a nonzero pairing class along some directions, and that class is fixed by data the framework supplies: the vacuum field selected by the curvature (the radial pair of Fact \ref{c:pairing.imported-schwarzschild-connection-curvature}, its radial member by the sign of the square, and a null direction in it by the flag bit) and the connection. A pairing generator exists and is derived, not premised, up to the flag bit.

\begin{itemize}
\item[K2.1] (nothing pairs) If the pairing class vanishes for every direction at every radius, the connection preserves the vacuum field and the claim is dead.
\item[K2.2] (frame artefact) If the class is nonzero but is not the transverse derivative $\ip{\nabla_X \ell}{m}$ of the vacuum's null direction, it is a property of the frame and not of the vacuum field, and the result is a reading, not a derivation. Test: the identity must hold in $\Q(F)$, and the vanishing locus must be invariant under $\ell \mapsto \lambda\ell$.
\item[K2.3] (consistency with the singularity paper) If the class along $\partial_t$ or $\partial_r$ is nonzero, the diagonality of that paper's supplied class \cite[Prop.~7.31, Thm.~C.8]{ButlerT5} is contradicted at the pre-ascent level; this would kill that paper's statement, not the claim, and would be reported as such.
\item[K2.4] (intra-fibre only) If every class found acts within the fibre at one point, and no component of the connection pairs the ladder at $x$ with the ladder at $y \neq x$, the cross-point clause of (L$_E$) is untouched and the claim closes only in the intra-fibre sense. Parallel transport along a path identifies the fibre at $x$ with the fibre at $y$, so ``transport, then compare'' is the intra-fibre result along a path; the cross-point object is a generator in $\so(V_{x} \oplus V_{y})$ for the completion's independent ladders, which no connection supplies. This condition bounds the claim and was expected to trigger.
\end{itemize}
Predictions, as registered: the class vanishes along $\partial_t$ and $\partial_r$, is proportional to $\sqrt f/r$ in the transverse directions, vanishes at the horizon and changes sign across it (P2.1); it equals $\ip{\nabla_X\ell}{m}$ (P2.2); the mixed-plane curvature has ``nonzero pairing class in the $n$-$m$ planes and zero in the $\ell$-$m$ planes'', and the radial and angular planes have none (P2.3); in the flat conformal model the classes of a displacement $dx$ are $(\ip{dx}{\ell}, \ip{dx}{m}, 0)$ and vanish only along the ray (P2.4).

\phantomsection\label{c:note.pairing.setting-two-claims-kill-conditions-fixed-3}
\textbf{Claim G1 (the mismatch angle).} The Hawking mismatch angle $\tan\theta = e^{-\omega/2T}$ of the singularity paper \cite[Lem.~7.2, Thm.~7.27]{ButlerT5} is the pairing class of a transport of the framework's connection: some path or loop has a holonomy with that pairing angle.

\begin{itemize}
\item[K1.1] (frame mismatch is Cartan) If the Schwarzschild static frame and the regular infalling frame differ at a point by an element of the stabiliser, their vacua coincide and the angle is not a frame mismatch at a point.
\item[K1.2] (holonomy carries no mode weight) If the class of every holonomy is a function of the geometry alone, the point and the loop, while the Hawking angle is a function of the mode weight $\omega$ and the period $\beta$, a one-parameter family in $\omega$, then no single holonomy reproduces the family.
\item[K1.3] (the sign of the square) If the Hawking pairing needs the Euclidean rotor, imaginary rapidity in the boost plane, which the singularity paper shows is not a member of the curvature-selected pair \cite[Prop.~7.21(3)]{ButlerT5}, while every holonomy is a real rotor, the claim is dead on this route too. This is that paper's result and is cited, not re-derived.
\end{itemize}
Predictions: K1.1 triggers (P1.1); K1.2 triggers (P1.2); the claim is not closed and the regularity premise remains the import.

\phantomsection\label{c:note.pairing.setting-two-claims-kill-conditions-fixed-4}
\textbf{Verdict rules.} R2 counts as closed in the intra-fibre sense if none of K2.1 to K2.3 triggers, and in the cross-point sense only if K2.4 also fails to trigger. G1 counts as closed only if none of K1.1 to K1.3 triggers. Prior art is assigned after the computation by the series' rule: where the identification of the pairing class is a known dictionary, the result is an import with a framework reading and is labelled so.

\section{The vacuum's stabiliser, and the pairing functional}

\begin{proposition}[The stabiliser is the parabolic algebra of the null direction]\label{c:pairing.stabiliser-parabolic-algebra-null-direction}
Let $W = \operatorname{span}\{\ell, m\} \subset \C^{4}$ and define, on $\so(4,\C)$, the linear functional
\[
\pv(B) := \ip{B\cdot\ell}{m} .
\]
Then a bivector $B$ stabilises $W$ (that is, $B\cdot W \subseteq W$, which is the condition for the spin action of $B$ to fix the vacuum ray) if and only if $\pv(B) = 0$. On the basis, $\pv(e_{01}) = \pv(e_{23}) = 0$, $\pv(e_{02}) = 1$, $\pv(e_{03}) = i$, $\pv(e_{12}) = -1$, $\pv(e_{13}) = -i$; the stabiliser has complex dimension five, and its real points form the four-dimensional parabolic algebra of the null direction $\ell$,
\[
\mathfrak{p}(\ell) = \operatorname{span}_{\R}\{\,e_{01},\ e_{23},\ \ell e_{2},\ \ell e_{3}\,\},
\]
the boost in the plane of $\ell$, the rotation about its axis, and the two null rotations fixing $\ell$. The null rotations about $n$ move the vacuum: $\pv(n e_{2}) = 2$, $\pv(n e_{3}) = 2i$. The class of $B$ is the $\mb$-component of $B\cdot\ell$: $\pv(B) = 0$ exactly when $B\cdot\ell \in \operatorname{span}\{\ell, m\}$.
\end{proposition}
\begin{proof}
$W$ is maximal isotropic, so $W^{\perp} = W$ and a vector $u$ lies in $W$ if and only if $\ip{u}{\ell} = \ip{u}{m} = 0$. For $B \in \so(4,\C)$ the action is skew, $\ip{B\cdot u}{v} = -\ip{u}{B\cdot v}$ (\ref{c:std-algebra.bivector-lie-algebra}), so $\ip{B\cdot\ell}{\ell} = \ip{B\cdot m}{m} = 0$ automatically and $\ip{B\cdot m}{\ell} = -\pv(B)$. Hence $B\cdot\ell \in W$ and $B\cdot m \in W$ together reduce to the single condition $\pv(B) = 0$, and the stabiliser is the kernel of a nonzero functional on a six-dimensional space. The values on the basis are the commutators $\tfrac12[e_{ab}, \ell]$ paired with $m$; for instance $\tfrac12[e_{02}, e_{0} + e_{1}] = -e_{2}$ and $\ip{-e_{2}}{e_{2} + ie_{3}} = 1$. The four real elements listed have $\pv = 0$ (the null rotations because $\tfrac12[\ell e_{2}, \ell] = 0$, $\ell$ being null and anticommuting with $e_{2}$ (\ref{c:std-algebra.orthogonal-anticommute})) and are linearly independent; the two real equations $\operatorname{Re}\pv = \operatorname{Im}\pv = 0$ are independent (the script checks the rank), so their real solutions form a space of dimension exactly four, which the four elements span. That the spin action of $B$ fixes the ray of the vacuum exactly when $B$ stabilises its annihilator $W$ is the definition of the annihilator, Fact \ref{c:pairing.imported-two-readings-capacities-vacuum-null}. The script verifies every value, the equivalence ``$B\cdot W \subseteq W$ if and only if $\pv(B) = 0$'' on twenty random complex bivectors, the independence, and the two null rotations about $n$.

\end{proof}

\begin{remark}[Reading]\label{c:note.pairing.reading}
By Fact \ref{c:pairing.imported-null-flags} the vacuum ray is the null direction and the orbit of vacua is the celestial sphere; Proposition \ref{c:pairing.stabiliser-parabolic-algebra-null-direction} is the statement, in the series' vocabulary, that the stabiliser of a point of the celestial sphere in the Lorentz algebra is the parabolic subalgebra of that null direction, and that the tangent of the sphere at $\ell$ is $\ell^{\perp}/\ell \otimes \C$, the transverse plane. What is the note's own is the functional $\pv$ in the Clifford conventions of the series and the next proposition, which says what $\pv$ is on the Fock module.
\end{remark}

\begin{proposition}[The pairing functional is the pair-creation amplitude]\label{c:pairing.pairing-functional-pair-creation-amplitude}
On the two-mode Fock module of the section ``Setting, and what is imported'', with $\ket{12} = a_{1}^{\dagger}a_{2}^{\dagger}\ket{0}$, every bivector $B$ of $\so(4,\C)$ satisfies
\[
\bra{12}\,B\,\ket{0} \;=\; -\,\pv(B) .
\]
On the even sector $\operatorname{span}\{\ket{0}, \ket{12}\}$ at $\nu = 2$ a bivector is a $2 \times 2$ matrix: the hops $a_{1}^{\dagger}a_{2}$, $a_{2}^{\dagger}a_{1}$ kill both states, the annihilation $a_{1}a_{2}$ kills $\ket{0}$ and sends $\ket{12}$ to $-\ket{0}$, the creation $a_{1}^{\dagger}a_{2}^{\dagger}$ sends $\ket{0}$ to $\ket{12}$, and the Cartan elements are diagonal. On the vacuum, therefore, a bivector acts through its Cartan scalar and its pairing class alone.
\end{proposition}
\begin{proof}
Both sides are linear in $B$, so it suffices to check the six basis bivectors, which the script does by forming the matrices of $e_{0}, \dots, e_{3}$ on the module, the Fock representation of \ref{c:std-algebra.fock-representation} (verified to satisfy the Clifford relations), multiplying, and reading off the $(\ket{12}, \ket{0})$ entry: $0, -1, -i, 1, i, 0$ against $\pv = 0, 1, i, -1, -i, 0$. The sector statements are the definition of the operators: a hop conserves the occupation number and $\ket{0}$, $\ket{12}$ have numbers $0$ and $2$ with no other states at those numbers; $a_{1}a_{2}\ket{0} = 0$, $a_{1}a_{2}\ket{12} = -\ket{0}$ and $a_{1}^{\dagger}a_{2}^{\dagger}\ket{0} = \ket{12}$ are computed in the script. This is Fact \ref{c:pairing.imported-stabiliser-tangent}'s decomposition read at $\nu = 2$, where the tangent is one-dimensional and $\pv$ is its coordinate.

\end{proof}

\begin{corollary}[The singularity paper's witness is a rotation of the null direction]\label{c:pairing.singularity-papers-witness-rotation-null-direction}
The rotor $R = \tfrac35 + \tfrac45\, e_{12}$ acts on the Fock module by $R\ket{0} = \tfrac35\ket{0} + \tfrac45\ket{12}$, with no odd component, and on vectors by $\ell \mapsto R\ell\tilde R = e_{0} - \tfrac{7}{25}e_{1} + \tfrac{24}{25}e_{2}$, a null vector not proportional to $\ell$ ($\ip{R\ell\tilde R}{\ell} = \tfrac{32}{25}$). The state $R\ket{0}$ is annihilated by $R\ell\tilde R$ and $Rm\tilde R$, so it is the Witt vacuum of the rotated decomposition. The necessity witness of Fact \ref{c:pairing.imported-stabiliser-tangent} is therefore the spin action of a real rotation of the frame about the $e_{3}$ axis, by the double angle with $(\cos, \sin) = (-\tfrac{7}{25}, \tfrac{24}{25})$: a mismatch of two Witt decompositions at a point is half the angle between two null directions, the half being Fact \ref{c:pairing.imported-null-flags}'s and the instance the one computed. By contrast the boost rotor $\tfrac53 + \tfrac43\,e_{01}$ of the singularity paper's Step 4a acts on $\ket{0}$ as the scalar $\tfrac13$ and fixes the ray.
\end{corollary}
\begin{proof}
Direct computation, all of it in the script: $R\tilde R = 1$; the column of $R$ on $\ket{0}$; the conjugate $R\ell\tilde R$ and its norm; the annihilation of $R\ket{0}$ by the two rotated null vectors; and the column of the boost rotor. The half angle is Fact \ref{c:pairing.imported-null-flags}; the real rotation $e_{12}$ acts on the vacuum purely through its pairing class by Proposition \ref{c:pairing.pairing-functional-pair-creation-amplitude}, since $\pv(e_{12}) = -1$, the hop and annihilation parts of $e_{12}$ kill $\ket{0}$, and $e_{12}$ has no Cartan part.

\end{proof}

\section{Transport moves the vacuum by the transverse derivative of its null direction}

\phantomsection\label{c:note.pairing.setting-transport-moves-vacuum-transverse-derivative}
Let a vacuum field be given on a region: a null direction field $\ell(x)$ with a transverse complex null vector $m(x)$, and let $\nabla$ be the Levi-Civita connection. In an orthonormal frame with $\ell = e_{0} + e_{1}$ and $m = e_{2} + ie_{3}$ of constant components, $\nabla_{X} e_{b} = \omega^{a}{}_{b}(X)\,e_{a}$, and the connection acts on the Clifford fibre by the bivector
\[
\Omega(X) \;=\; \sum_{a<b}\omega^{ab}(X)\,e_{a}e_{b}, \qquad \omega^{ab} = \omega^{a}{}_{c}\,\eta^{cb},
\]
in the sense that $\nabla_{X} v = \tfrac12[\Omega(X), v]$ for every vector $v$ of constant frame components. The spin transport of the vacuum along $X$ is the spin action of the same $\Omega(X)$.

\begin{definition}[The pairing class of the transport]\label{c:obj.pairing.pairing-class-transport}
The pairing class of the transport along $X$, relative to the vacuum field, is $\pv(\Omega(X)) = \ip{\Omega(X)\cdot\ell}{m}$. The transport preserves the vacuum ray along $X$ if and only if $\pv(\Omega(X)) = 0$ (Proposition \ref{c:pairing.stabiliser-parabolic-algebra-null-direction}), and by Proposition \ref{c:pairing.pairing-functional-pair-creation-amplitude} the spin transport $\exp(s\,\Omega(X))$ carries the vacuum, to first order in $s$ and up to a scalar, to $\ket{0} - s\,\pv(\Omega(X))\,\ket{12}$: the class is the Thouless amplitude of the mismatch. (Comparing the vacuum field against the transported frame instead reverses the sign; nothing below depends on the convention.)
\end{definition}

\begin{proposition}[The pairing class is the transverse derivative of the null direction]\label{c:pairing.pairing-class-transverse-derivative-null-direction}
For every direction $X$,
\[
\pv(\Omega(X)) \;=\; \ip{\nabla_{X}\ell}{m},
\]
and under a rescaling $\ell \mapsto \lambda\ell$ of the null direction by a function the class scales by $\lambda$, so its vanishing locus is a property of the null direction field and not of its normalisation. For a real direction $X$ the vacuum is parallel along $X$ exactly when the null direction is, $\nabla_{X}\ell \propto \ell$.
\end{proposition}
\begin{proof}
$\nabla_{X}\ell = \tfrac12[\Omega(X), \ell] = \Omega(X)\cdot\ell$ by the defining property of $\Omega$, so the identity is the definition of $\pv$ once $\Omega$ reproduces $\nabla$, which is what the script verifies for both Schwarzschild frames. For the rescaling, $\nabla_{X}(\lambda\ell) = \lambda\nabla_{X}\ell + (X\lambda)\ell$ and $\ip{\ell}{m} = 0$; the script models the derivative term, which is nonzero along $e_{1}$, and confirms the scaling in every direction with $\lambda = F/(1+F)$. For real $X$, $\pv(\Omega(X)) = 0$ says $\nabla_{X}\ell \in W$ (by \ref{c:pairing.stabiliser-parabolic-algebra-null-direction}), and $\nabla_{X}\ell$ is real, so it is a real vector of $W = \operatorname{span}\{\ell, m\}$, hence a multiple of $\ell$. (For a complex direction the equivalence fails: in Schwarzschild $\nabla_{m}\ell \propto m$ with $\pv(\Omega(m)) = 0$.)

\end{proof}

\begin{remark}[Reading, and prior art]\label{c:note.pairing.reading-prior-art}
Proposition \ref{c:pairing.pairing-class-transverse-derivative-null-direction} is elementary once Proposition \ref{c:pairing.stabiliser-parabolic-algebra-null-direction} is in hand, and its content is a dictionary: by Fact \ref{c:pairing.imported-spin-coefficients-null-tetrad} the components of the pairing class on the null tetrad are, up to normalisation and sign convention, the spin coefficients $\kappa, \tau, \sigma, \rho$, so the pairing generator of the framework's vacuum is the transverse deformation of the vacuum's null congruence, the optical scalars of Sachs and Newman-Penrose: the vacuum is parallel along its own integral curve if $\ell$ is geodesic ($\kappa = 0$), along $n$ if $\tau = 0$, and along every transverse real direction if the congruence has neither expansion, twist nor shear ($\rho = \sigma = 0$). Under the series' rule this is an import with a framework reading. The reading is this. The stabiliser part of the connection, eight of the twelve spin coefficients, rescales $\ell$, rotates $m$ in its plane and tilts $n$ and $\mb$ by null rotations about $\ell$; it is what the singularity paper's hops, annihilations and Cartan elements do, and it cannot move the vacuum, and so cannot create entanglement from it (off the vacuum a hop does entangle, as that paper's own script exhibits \cite[Rem.~13.5]{ButlerT5}). The pairing part, the other four, tilts $\ell$ and is what the framework had said it did not supply. It supplies it wherever the congruence of its chosen null direction is not parallel.
\end{remark}

\section{Schwarzschild, exactly}

\phantomsection\label{c:note.pairing.setting-schwarzschild-exactly}
The framework's Schwarzschild vacuum is the radial one: the extreme curvature eigenvalue distinguishes the pair of the radial plane and its dual, the sign of the square selects the radial member (Fact \ref{c:pairing.imported-schwarzschild-connection-curvature}), the flag bit picks the outgoing or the ingoing null direction in it, and the chirality bit picks $m$. The computations below take $\ell$ outgoing, and Proposition \ref{c:pairing.regular-frame-class-vanishes-horizon-changes} records what the ingoing choice gives instead; they are exact in the field $\Q(F)$, $F = \sqrt f$, $r = 2M/(1 - F^{2})$, at $M = \tfrac12$ (the singularity paper's device for keeping every coefficient rational), and in $\Q(\beta)$, $r = 2M/\beta^{2}$, for the regular frame.

\begin{proposition}[The static frame: only the expansion survives]\label{c:pairing.static-frame-only-expansion-survives}
For the Schwarzschild gravitational field on flat spacetime in the static gauge (Definition \emph{Gauge Theory Gravity: the gravitational field on flat spacetime}), with the connection of Fact \ref{c:pairing.imported-schwarzschild-connection-curvature} read along the static orthonormal frame of physical vectors $\{e_a\}$, and $\Omega(e_c)$ written for the rotation-gauge field along the physical vector $e_c$ (the bivector $\Omega(\hb(e_c))$ of item (b) of the definition), the bivectors $\Omega(e_{c})$ reproduce the covariant derivatives $(e_{c} \cdot D)\, e_{b} = \omega^{a}{}_{b}(e_{c})e_{a}$ for all $b, c$, and the pairing classes are
\[
\pv(\Omega(e_{0})) = \pv(\Omega(e_{1})) = 0, \qquad \pv(\Omega(e_{2})) = -\frac{\sqrt f}{r}, \qquad \pv(\Omega(e_{3})) = -\,i\,\frac{\sqrt f}{r},
\]
so that on the null tetrad $\pv(\Omega(\ell)) = \pv(\Omega(n)) = \pv(\Omega(m)) = 0$ and $\pv(\Omega(\mb)) = -2\sqrt f/r$: the congruence is geodesic, $\tau = 0$, shear-free, and its only optical scalar is the transverse expansion $2\sqrt f/r$. Along its own direction $(\ell \cdot D)\ell = (f'/2\sqrt f)\,\ell$, a rescaling. In particular $\Omega(e_{0})$ is a multiple of $e_{01}$ and $\Omega(e_{1}) = 0$: the transports along the $t$-line and the $r$-line lie in the stabiliser. Every class vanishes at $F = 0$, and at $r = 25/16$ ($F = \tfrac35$) the transverse class is $-\tfrac{48}{125}$.
\end{proposition}
\begin{proof}
Read in Gauge Theory Gravity, the connection forms and structure equations below are those of the rotation-gauge field in the orthonormal coframe of the static gauge, and $\nabla$ is the covariant derivative $D$ (Fact \ref{c:synthesis.imported-gtg-dictionary}, clauses (i) and (ii)). The connection forms of \ref{c:pairing.imported-schwarzschild-connection-curvature} give $\omega^{a}{}_{b}(e_{c})$ as elements of $\Q(F)$: $f' = 1/r^{2} = (1 - F^{2})^{2}$ (an identity the script checks), so $f'/2\sqrt f = (1 - F^{2})^{2}/2F$ and $\sqrt f/r = F(1 - F^{2})$. The script verifies these forms against the first structure equations on the $(t, r, \vartheta)$ block, with $f'$ obtained by differentiation, which with the antisymmetry of $\omega_{ab}$ fixes the Levi-Civita connection; forms $\Omega(e_{c})$, verifies it reproduces every frame derivative as an identity of rational functions, and computes the classes. The transverse values are $\ip{\omega^{2}{}_{1}(e_{2})\,e_{2}}{m} = (\sqrt f/r)\ip{e_{2}}{e_{2} + ie_{3}} = -\sqrt f/r$ and likewise for $e_{3}$; along $e_{0}$ the derivative of $\ell$ is $(f'/2\sqrt f)(e_{1} + e_{0})$, proportional to $\ell$, and along $e_{1}$ it vanishes. The tetrad values follow by linearity, $\ip{\mb}{m} = -2$. Read through Fact \ref{c:pairing.imported-spin-coefficients-null-tetrad}, the vanishing classes on $\ell$, $n$ and $m$ are the vanishing spin coefficients $\kappa$, $\tau$ and $\sigma$: the congruence is geodesic, $\tau = 0$ and shear-free. The transverse expansion, the trace $\sum_{c = 2,3}\eta^{cc}\ip{e_{c}}{\nabla_{e_{c}}\ell}$, is $2\sqrt f/r$ and equals $-\pv(\Omega(\mb))$. The script also checks that the block one-forms and the frame matrices it uses for the classes are the same connection, evaluated on the frame vectors.

\end{proof}

\begin{remark}[Reading]\label{c:note.pairing.reading-2}
That the $t$-line and $r$-line transports lie in the stabiliser is a pre-ascent statement about the Dirac fibre. The singularity paper's diagonality is a statement about its supplied class on the boundary modes \cite[Prop.~7.31, Thm.~C.8]{ButlerT5}, a different space; the two agree in kind, the radial-plane boost being the Cartan element in both, and this note does not claim the one proves the other.
\end{remark}

\begin{proposition}[The regular frame: the class vanishes on the horizon and changes sign inside]\label{c:pairing.regular-frame-class-vanishes-horizon-changes}
For the Schwarzschild gravitational field on flat spacetime in the Painlev\'e-Gullstrand gauge $\hb = 1 - \beta P$ (Definition \emph{Gauge Theory Gravity: the gravitational field on flat spacetime}), take the orthonormal physical frame $\{e_a\}$ whose coframe is the set of components of the physical displacement $\hb^{-1}(dx)$: $\theta^{0} = dT$, $\theta^{1} = dr + \beta\,dT$, $\theta^{2} = r\,d\vartheta$, $\theta^{3} = r\sin\vartheta\,d\varphi$. The rotation-gauge field along this frame, written $\Omega(e_c)$ for $\Omega(\hb(e_c))$ (the bivector of item (b) of the definition along the physical vector $e_c$), has the connection forms $\omega^{0}{}_{1} = -\beta'\,\theta^{1}$, $\omega^{2}{}_{1} = \theta^{2}/r$, $\omega^{2}{}_{0} = -(\beta/r)\,\theta^{2}$, with $\hat3$ mirroring $\hat2$ and $\omega^{3}{}_{2} = (\cot\vartheta/r)\theta^{3}$, where $\beta' = d\beta/dr = -\beta^{3}/4M$. The pairing classes of the outgoing vacuum $\ell = e_{0} + e_{1}$ are
\[
\pv(\Omega(e_{0})) = \pv(\Omega(e_{1})) = 0, \qquad \pv(\Omega(e_{2})) = -\frac{1 - \beta}{r}, \qquad \pv(\Omega(e_{3})) = -\,i\,\frac{1 - \beta}{r},
\]
zero on the horizon locus $\beta = 1$, negative outside and positive inside; the transverse expansion is $2(1 - \beta)/r$, the expansion of the outgoing congruence, which vanishes on the horizon locus and is negative in the trapped region; and $(\ell \cdot D)\ell = (\beta^{3}/4M)\,\ell$. For the ingoing flag bit, the vacuum $\operatorname{span}\{n, m\}$ with $n = e_{0} - e_{1}$, the transverse classes are $+(1 + \beta)/r$ and $+i(1 + \beta)/r$: one sign throughout, $\tfrac{144}{125}$, $2$ and $12$ at $\beta = \tfrac45, 1, 2$, and no horizon zero, the ingoing rays converging everywhere. The two frames agree covariantly: $\pv_{\mathrm{PG}} = \tfrac{F}{1 + \beta}\,\pv_{\mathrm{static}}$ in both transverse directions, the factor being the rescaling $\ell_{\mathrm{PG}} = (F/(1 + \beta))\,\ell_{\mathrm{static}}$; at $(F, \beta) = (\tfrac35, \tfrac45)$ the values are $-\tfrac{16}{125} = \tfrac13\cdot(-\tfrac{48}{125})$. The field strength $\mathcal{R}(B)$ read in the regular frame has the same coefficients as functions of $r$, and the class of the loop $(n, e_{2})$ is $2M/r^{3}$ in this frame too, equal to $1$ on the horizon locus at $M = \tfrac12$: the curvature classes do not vanish where the connection classes do.
\end{proposition}
\begin{proof}
Read in Gauge Theory Gravity, the structure equations below are the vanishing torsion of the rotation-gauge field written in the orthonormal coframe of the Painlev\'e-Gullstrand gauge, the components of $\hb^{-1}(dx)$, and the Levi-Civita derivative is $D$ (Fact \ref{c:synthesis.imported-gtg-dictionary}, clauses (i) and (ii)). The coframe is orthonormal for the Painlev\'e-Gullstrand metric of Fact \ref{c:pairing.imported-schwarzschild-connection-curvature}. The connection is found from the first structure equations $d\theta^{a} + \omega^{a}{}_{b}\wedge\theta^{b} = 0$: $d\theta^{0} = 0$ is consistent with $\omega^{0}{}_{1} \propto \theta^{1}$; $d\theta^{1} = \beta'\,dr\wedge dT = \beta'\,\theta^{1}\wedge\theta^{0}$ fixes the coefficient; $d\theta^{2} = dr\wedge d\vartheta = (\theta^{1} - \beta\theta^{0})\wedge\theta^{2}/r$ gives $\omega^{2}{}_{1}$ and $\omega^{2}{}_{0}$; the $\hat3$ forms follow by spherical symmetry as in Fact \ref{c:pairing.imported-schwarzschild-connection-curvature}. The script verifies that these forms satisfy all three structure equations exactly on the $(T, r, \vartheta)$ block, with $d$ acting on functions of $r$ through $d/dr = \beta'\,d/d\beta$ in $\Q(\beta)$, which with antisymmetry determines the Levi-Civita connection; it then forms the bivectors, checks them against the frame derivatives, and computes the classes: $\nabla_{e_{2}}\ell = (\omega^{2}{}_{0} + \omega^{2}{}_{1})(e_{2})\,e_{2} = ((1 - \beta)/r)\,e_{2}$. The sign change is the sign of $1 - \beta$, and by Fact \ref{c:pairing.imported-spin-coefficients-null-tetrad} the class on $\mb$, $-2(1 - \beta)/r$, is the spin coefficient $\rho$, minus the expansion of the outgoing congruence; the evaluations at $\beta = \tfrac45, 1, 2$ give $-\tfrac{16}{125}, 0, 4$ at $M = \tfrac12$. For the ingoing bit, $\nabla_{e_{2}}n = (\omega^{2}{}_{0} - \omega^{2}{}_{1})(e_{2})\,e_{2} = -((1 + \beta)/r)\,e_{2}$, paired with $m$; the script computes the classes of $\operatorname{span}\{n, m\}$ and the three values. For the covariance: the static and Painlev\'e-Gullstrand frames share $e_{2}, e_{3}$ and their null directions are proportional, Proposition \ref{c:pairing.two-frames-share-vacuum-selected-transport}(i), with $\ell_{\mathrm{PG}} = ((1 - \beta)/F)\,\ell_{\mathrm{static}} = (F/(1 + \beta))\,\ell_{\mathrm{static}}$ on $F^{2} + \beta^{2} = 1$, so Proposition \ref{c:pairing.pairing-class-transverse-derivative-null-direction}'s scaling law applies with a $\lambda$ depending on $r$ alone; the script checks the relation at five rational points of the circle, $F = (1 - u^{2})/(1 + u^{2})$, $\beta = 2u/(1 + u^{2})$. The curvature is derived in this frame by the second structure equations, as in the static frame, and its $(n, e_{2})$ class is compared with the static one at three rational radii; it is invariant under the radial boost because $\ell$ and $n$ scale inversely.

\end{proof}

\begin{remark}[What the horizon is, in this reading]\label{c:note.pairing.what-horizon-reading}
The horizon is where the pairing generator of the outgoing vacuum vanishes in every direction. Along $\ell$ and along $n$ it vanishes at every radius, the congruence being geodesic and $\tau = 0$ throughout Schwarzschild; what singles out the horizon is the transverse class, which vanishes there and nowhere else because the expansion of the horizon generators is zero, a Killing-horizon fact (Fact \ref{c:pairing.imported-spin-coefficients-null-tetrad}) that Proposition \ref{c:pairing.regular-frame-class-vanishes-horizon-changes} exhibits exactly. In the trapped region the class has the opposite sign: transport tilts the outgoing null direction the other way. None of this uses the ascent; it is the pre-ascent geometry of the null direction the descent selects, and it is where the singularity paper's Step 4d premise is calibrated, at the fixed locus of the boost \cite[Prop.~7.25]{ButlerT5}. The evidence for the horizon statement is the regular-frame computation of Proposition \ref{c:pairing.regular-frame-class-vanishes-horizon-changes}, not the static one: the static frame degenerates on the horizon and the rescaling between the two frames vanishes there, so Proposition \ref{c:pairing.pairing-class-transverse-derivative-null-direction}'s invariance does not carry the static result across; that the static class $-\sqrt f/r$ also vanishes is consistent, not evidential.
\end{remark}

\begin{proposition}[Curvature: which loops move the vacuum]\label{c:pairing.curvature-which-loops-move-vacuum}
For the Schwarzschild gravitational field on flat spacetime (Definition \emph{Gauge Theory Gravity: the gravitational field on flat spacetime}) and the loop spanned by physical vectors $e_{c}, e_{d}$, let $R(e_{c}, e_{d}) = R(\hb(e_{c} \wedge e_{d})) = -\mathcal{R}(e_{c} \wedge e_{d})$ be the curvature bivector, the covariant field strength on the physical plane $e_{c} \wedge e_{d}$ in the sign of \cite{LDG1998}, so that the holonomy of a small loop of area $\epsilon$ is $\exp(\epsilon R)$ to first order. In the static orthonormal frame of physical vectors the pairing classes are
\begin{gather*}
\pv(R(e_{0},e_{1})) = \pv(R(e_{2},e_{3})) = 0, \\
\pv(R(e_{0},e_{2})) = \frac{M}{r^{3}}, \qquad \pv(R(e_{1},e_{2})) = -\frac{M}{r^{3}}, \qquad \pv(R(e_{0},e_{3})) = i\frac{M}{r^{3}}, \qquad \pv(R(e_{1},e_{3})) = -i\frac{M}{r^{3}} .
\end{gather*}
Hence $R(\ell, e_{2}) = (M/r^{3})\,\ell e_{2}$ is a null rotation about $\ell$ with zero class, while $R(n, e_{2})$ has class $2M/r^{3}$; on the complex loops, $(\ell, m)$, $(\ell, \mb)$, $(m, \mb)$ and $(n, m)$ have zero class and $(n, \mb)$ has class $4M/r^{3}$. The holonomy therefore moves the vacuum exactly for loops with an $n \wedge \mb$ component, and for no loop lying in the null hypersurface generated by $\ell$.
\end{proposition}
\begin{proof}
Read in Gauge Theory Gravity, the connection and curvature forms below are the rotation-gauge field and the field strength of the Schwarzschild field in the static orthonormal frame, and the bivector they give on the plane $e_{c} \wedge e_{d}$ is $R(\hb(e_{c} \wedge e_{d})) = -\mathcal{R}(e_{c} \wedge e_{d})$ (Fact \ref{c:synthesis.imported-gtg-dictionary}, clauses (ii) and (iv)). The script derives the curvature two-forms on the $(t, r, \vartheta)$ block from the connection by the second structure equations $R^{a}{}_{b} = d\omega^{a}{}_{b} + \omega^{a}{}_{c}\wedge\omega^{c}{}_{b}$, with $f'$ and $f''$ obtained by differentiation in $\Q(F)$, and finds $R^{0}{}_{1} = -(f''/2)\,\theta^{0}\wedge\theta^{1}$, $R^{0}{}_{2} = R^{1}{}_{2} = -(f'/2r)\,\theta^{0}\wedge\theta^{2}, \theta^{1}\wedge\theta^{2}$, with $-f''/2 = 2M/r^{3}$ and $-f'/2r = -M/r^{3}$; the $\hat3$ planes mirror $\hat2$ and the angular coefficient $(1 - f)/r^{2}$ is taken from Fact \ref{c:pairing.imported-schwarzschild-connection-curvature}. It then lowers the index, forms the bivectors and computes the six classes; the loop statements are bilinearity, and the Weyl reading of them is Lemma \ref{c:pairing.loop-classes-are-weyl-components}. That $R(\ell, e_{2})$ is a null rotation about $\ell$ is the identity $(M/r^{3})(e_{02} + e_{12}) = (M/r^{3})\,\ell e_{2}$, checked as multivectors.

\end{proof}

\begin{fact}[Imported: the Weyl scalars, and the Petrov type of Schwarzschild]\label{c:pairing.imported-weyl-scalars-petrov-type-schwarzschild}
For a gravitational field $(\hb, \Omega)$ on flat spacetime (Definition \emph{Gauge Theory Gravity: the gravitational field on flat spacetime}) and a null tetrad $(\ell, n, m, \mb)$ of physical vectors, the Weyl map $\mathcal{W}$ of the field strength, read as the Weyl tensor $C$ of the derived metric on the tetrad's bivectors $\ell \wedge m$, $\ell \wedge n$ and $n \wedge \mb$, has five complex components, $\Psi_{0} = C(\ell, m, \ell, m)$, $\Psi_{1} = C(\ell, n, \ell, m)$, $\Psi_{2} = C(\ell, m, \mb, n)$, $\Psi_{3} = C(\ell, n, \mb, n)$, $\Psi_{4} = C(n, \mb, n, \mb)$, up to normalisation and sign convention; $\ell$ is a principal null direction if and only if $\Psi_{0} = 0$ and a repeated one if and only if $\Psi_{0} = \Psi_{1} = 0$. The Schwarzschild field is of Petrov type D with the outgoing and ingoing radial null directions as its repeated principal directions, $\Psi_{2} = -M/r^{3}$ in the unit-normalised tetrad and all other $\Psi$ zero. For a vacuum field, $\mathcal{R}(a) = 0$ (Ricci-flat), the curvature operator is its Weyl map, $\mathcal{R} = \mathcal{W}$. \emph{Provenance:} \cite[Ch.~8]{PenroseRindler2}; \cite{NewmanPenrose1962}.
\emph{Source:} R. Penrose, W. Rindler, Spinors and Space-Time, Vol. 2: Spinor and Twistor Methods in Space-Time Geometry, Cambridge University Press (1986)., Ch.~8; E. Newman, R. Penrose, An approach to gravitational radiation by a method of spin coefficients, J. Math. Phys. 3 (1962) 566--578..
\end{fact}

\begin{lemma}[The loop classes are Weyl components]\label{c:pairing.loop-classes-are-weyl-components}
For a vacuum gravitational field on flat spacetime, $\mathcal{R}(a) = 0$ (Ricci-flat; Definition \emph{Gauge Theory Gravity: the gravitational field on flat spacetime}), let $R(u, v) = R(\hb(u \wedge v)) = -\mathcal{R}(u \wedge v)$ be the curvature bivector of the loop $(u, v)$ of physical vectors, the covariant field strength on the plane $u \wedge v$, and write
\[
C(u, v, w, z) := \ip{R(u, v)\cdot w}{z},
\]
with the flat product of physical vectors, so that the class of the loop $(u, v)$ is $C(u, v, \ell, m)$. The six loops of the null tetrad then have classes
\[
(\ell, m) : \Psi_{0}, \quad (\ell, \mb) : 0, \quad (m, \mb) : -\Psi_{1}, \quad (\ell, n) : \Psi_{1}, \quad (n, m) : 0, \quad (n, \mb) : -\Psi_{2},
\]
up to the tetrad's normalisation. Hence every loop in the null cone of $\ell$ has zero class if and only if $\Psi_{0} = \Psi_{1} = 0$, that is, if and only if $\ell$ is a repeated principal null direction; the radial loop then has zero class as well, and only $\Psi_{2}$ remains.
\end{lemma}
\begin{proof}
Read in Gauge Theory Gravity, $C$ below is the Weyl map of the vacuum field strength evaluated on physical vectors with the flat product, and its symmetries and tracelessness are those of $\mathcal{W}$ (Fact \ref{c:synthesis.imported-gtg-dictionary}, clauses (iv) and (v)). The loops $(\ell, m)$ and $(\ell, n)$ give the definitions of $\Psi_{0}$ and $\Psi_{1}$ (\ref{c:pairing.imported-weyl-scalars-petrov-type-schwarzschild}). The loop $(n, \mb)$ gives $C(n, \mb, \ell, m) = C(\ell, m, n, \mb) = -\Psi_{2}$, by the pair symmetry and the antisymmetry in the last pair. The remaining three use the tracelessness of the Weyl tensor, $\eta^{ab}C(e_{a}, u, e_{b}, v) = 0$, in the null tetrad, where the trace is $\tfrac12[C(\ell, u, n, v) + C(n, u, \ell, v)] - \tfrac12[C(m, u, \mb, v) + C(\mb, u, m, v)]$ up to the normalisation: with $u = v = \ell$ the two $\ell$-terms vanish by antisymmetry and the trace forces $C(m, \ell, \mb, \ell) = -C(\mb, \ell, m, \ell)$, both equal to $C(\ell, m, \ell, \mb)$ by the symmetries, hence zero; with $(u, v) = (m, \ell)$ it gives $C(\ell, m, n, \ell) = C(\mb, m, m, \ell)$, whose left side is $-\Psi_{1}$ and whose right side is $C(m, \mb, \ell, m)$; with $u = v = m$ it gives $2C(\ell, m, n, m) = 0$. The script verifies the six identities on the derived Schwarzschild curvature twice: in the adapted tetrad, where it also verifies Ricci-flatness on all sixteen frame pairs, $\Psi_{0} = \Psi_{1} = 0$ and $\Psi_{2} = -4M/r^{3}$ (the factor $4$ from $\ip{\ell}{n} = 2$ and $\ip{m}{\mb} = -2$), so that three of the six entries there read $0 = 0$; and in a generic rational null tetrad, rotated by a product of three rational rotors, where $\Psi_{0}$ and $\Psi_{1}$ are nonzero and all six entries are exercised as identities in $\Q(F)$.

\end{proof}

\begin{remark}[Reading]\label{c:note.pairing.reading-3}
Proposition \ref{c:pairing.curvature-which-loops-move-vacuum} is the component form of Fact \ref{c:pairing.imported-weyl-scalars-petrov-type-schwarzschild}'s $\Psi_{2}$ through Lemma \ref{c:pairing.loop-classes-are-weyl-components}: the loops that move the framework's Schwarzschild vacuum are exactly those carrying $\Psi_{2}$, and the null cone of $\ell$ carries none because $\ell$ is a repeated principal null direction. The framework's radial vacuum is therefore not only the one selected by the extreme curvature eigenvalue and the sign of the square; it is the one whose null cone the holonomy leaves fixed. This is an import with a reading.
\end{remark}

\section{The ascended reading: the point mode}

\phantomsection\label{c:note.pairing.setting-ascended-reading-point-mode}
In the ascended algebra the complex capacity is three, and the third mode is the point itself: the vacuum at $x$ is a twistor whose isotropic three-space contains $X(x)$ (Fact \ref{c:pairing.imported-flat-conformal-model-twistor-incidence}). The framework's vacuum field in the flat model is
\[
W(x) = \operatorname{span}\{\,X(x),\ \ell(x),\ m(x)\,\}, \qquad \ell(x) = \ell - \ip{x}{\ell}\,\ninf, \quad m(x) = m - \ip{x}{m}\,\ninf ,
\]
the lifts of the null tetrad orthogonal to $X(x)$, and the three pairing classes of a bivector $B$ relative to $W(x)$ are $\ip{B\cdot X}{\ell(x)}$, $\ip{B\cdot X}{m(x)}$, $\ip{B\cdot\ell(x)}{m(x)}$, one for each pair of modes.

\begin{proposition}[Displacement pairs the point mode; the vacuum is parallel along its ray only]\label{c:pairing.displacement-pairs-point-mode-vacuum-parallel}
$W(x)$ is totally isotropic of dimension three. The motion of the vacuum field under a displacement $dx = dx^{a}e_{a}$, the flat connection being the ambient derivative, is generated exactly by the translation bivector $\Omega(dx) = dx^{a}\,e_{a}\ninf$: $\tfrac12[e_{a}\ninf, X(x)] = \partial X/\partial x^{a} = e_{a} - x_{a}\ninf$, $\tfrac12[e_{a}\ninf, \ninf] = 0$, and the finite displacement $X(x + dx) - X(x)$ differs from $dx^{a}\,\partial_{a}X$ by $-\tfrac12 Q(dx)\ninf$ alone. The pairing classes of the transport are
\[
\bigl(\ip{dx}{\ell},\ \ip{dx}{m},\ 0\bigr):
\]
the displacement pairs the point mode with the boost mode by its $n$-component and with the angular mode by its $\mb$-component, and never pairs the boost mode with the angular mode. For a real displacement all three vanish if and only if $dx$ is a real multiple of $\ell$ (for a complex one the equivalence fails, $dx = m$ having all three classes zero): the vacuum is parallel along the point's own null ray and along no other real direction, which is the incidence relation of the conformal spinor space (twistor incidence), $X(x + s\ell) = X(x) + s\,\ell(x) \in W(x)$ for every $s$. The stabiliser of a $\nu = 3$ vacuum in $\so(6,\C)$ has dimension $12 = \nu^{2} + \nu(\nu-1)/2$, the three classes being independent functionals.
\end{proposition}
\begin{proof}
Isotropy: $X(x)$ is null by Fact \ref{c:pairing.imported-flat-conformal-model-twistor-incidence}, the lifts are null because $\ninf$ is null and orthogonal to $\ell$, $m$, and $\ip{X(x)}{\ell(x)} = \ip{x}{\ell} - \ip{x}{\ell}\ip{\nz}{\ninf} = 0$, likewise for $m$. The bivector identity is $\tfrac12[u\wedge v, w] = \ip{v}{w}u - \ip{u}{w}v$ (\ref{c:std-algebra.contraction-expansion}) with $u = e_{a}$, $v = \ninf$, $w = X(x)$ and $\ip{\ninf}{X} = 1$, $\ip{e_{a}}{X} = x_{a}$. The classes: $\ip{\Omega(dx)\cdot X}{\ell(x)} = \ip{dx^{a}(e_{a} - x_{a}\ninf)}{\ell(x)} = \ip{dx}{\ell}$ since $dx \perp \ninf$ and $\ninf \perp \ell(x)$; the same for $m$; and $\Omega(dx)\cdot\ell(x) = -\ip{dx}{\ell}\ninf$ is orthogonal to $m(x)$. A real $dx$ with $\ip{dx}{\ell} = \ip{dx}{m} = 0$ is orthogonal to $\ell, m, \mb$, hence to $\ell^{\perp}$, hence a multiple of $\ell$; the script checks the complex counterexample $dx = m$ and the ingoing $dx = \mb$ with classes $(0, -2, 0)$. Incidence is the expansion of $X(x + s\ell)$ with $Q(\ell) = 0$. The dimension is the rank of the three functionals on the fifteen basis bivectors of $\Cl(2,4)$. The script verifies each statement at four rational points and twenty random real displacements, the four named real ones and the two complex ones, the translation identity for each $a$, incidence at three values of $s$, and, for the dimension, that the common kernel of the three functionals has dimension twelve and coincides with the stabiliser on ten random kernel elements and on every basis bivector.

\end{proof}

\begin{remark}[What this does and does not say]\label{c:note.pairing.what-does-does-not-say}
Proposition \ref{c:pairing.displacement-pairs-point-mode-vacuum-parallel} is the incidence relation of the conformal spinor space (twistor incidence) in the vocabulary of pairing classes, an import with a reading: the point mode is annihilated by the point and created by $\ninf$, moving the point moves its vacuum, and the amount is the displacement's null components. It is stated for the flat model, where the laboratory paper states it, and not for the singular boundary, where the completion attaches an independent ladder pair at each point of the exceptional set and the tractor connection on the collar is not constructed in the series. Whether the transport along the singular boundary pairs the boundary mode at a point with the fibre modes there, as the flat model's translation pairs the point mode, is a successor problem, and this note does not claim it.
\end{remark}

\section{What this does to the two claims}

\begin{proposition}[The two frames share a vacuum, and the selected transport has no pairing class]\label{c:pairing.two-frames-share-vacuum-selected-transport}
For the Schwarzschild gravitational field on flat spacetime (Definition \emph{Gauge Theory Gravity: the gravitational field on flat spacetime}): (i) Outside the horizon locus, the static frame and the Painlev\'e-Gullstrand frame, two orthonormal frames of physical vectors, are related at every radius by $e_{0}^{\mathrm{PG}} = (e_{0} - \beta e_{1})/F$, $e_{1}^{\mathrm{PG}} = (e_{1} - \beta e_{0})/F$, a Lorentz transformation ($\Lambda^{\mathsf T}\eta\Lambda = \eta$), that is a rotation-gauge transformation, that fixes $e_{2}, e_{3}$ and scales $\ell$ by $(1 - \beta)/F = F/(1 + \beta)$: a boost in the plane of $\ell$, generated by $e_{01}$, which lies in $\mathfrak{p}(\ell)$. The two frames' vacua coincide; on the horizon locus itself the static frame does not exist, and the comparison there is carried by the regular frame alone, whose vacuum is the same object continued. (ii) The one transport the framework selects for a boundary mode is the flow along the $t$-line, whose holonomy is the boost of the completed plane \cite[Prop.~13.2(i)]{ButlerT5}; at the pre-ascent level its class is zero in both frames (Propositions \ref{c:pairing.static-frame-only-expansion-survives}, \ref{c:pairing.regular-frame-class-vanishes-horizon-changes}). (iii) At the singularity paper's sample point $e^{-\omega/T} = \tfrac{9}{16}$, if the Hawking states for the weights $\omega$ and $2\omega$ are written on the fibre's two modes as $\ket{0} + A\ket{12}$ with $A = \tfrac34$ and $A = \tfrac{9}{16}$, each is the vacuum of a decomposition obtained from $W$ by the rotor $1 + A e_{12}$, whose null directions are
\[
e_{0} + \tfrac{7}{25}e_{1} + \tfrac{24}{25}e_{2} \qquad\text{and}\qquad e_{0} + \tfrac{175}{337}e_{1} + \tfrac{288}{337}e_{2},
\]
distinct: a holonomy reproducing both would be one rotor per mode.
\end{proposition}
\begin{proof}
Read in Gauge Theory Gravity, the coordinate vector fields below are background directions, the frames are orthonormal sets of physical vectors, and the Lorentz matrix between them is a rotation-gauge transformation (Fact \ref{c:synthesis.imported-gtg-dictionary}, clause (i); Definition \emph{Gauge Theory Gravity: the gravitational field on flat spacetime}, item (e)). (i) With $T = t + h(r)$, $h' = \beta/f$, one has $\partial_{T} = \partial_{t}$ and $\partial_{r}|_{T} = \partial_{r}|_{t} - (\beta/f)\partial_{t}$, so $e_{0}^{\mathrm{PG}} = \partial_{T} - \beta\partial_{r}|_{T} = (1/f)\partial_{t} - \beta\partial_{r}|_{t} = (e_{0} - \beta e_{1})/F$ and $e_{1}^{\mathrm{PG}} = \partial_{r}|_{T} = (e_{1} - \beta e_{0})/F$ in terms of the static $e_{0} = \partial_{t}/F$, $e_{1} = F\partial_{r}$; the script checks $\Lambda^{\mathsf T}\eta\Lambda = \eta$, the images of $\ell$ and $e_{2}$, and the identity $(1 - \beta)/F = F/(1 + \beta)$ at three rational points of the circle. (ii) is the cited proposition (\ref{c:singularity.modular-flow-boundary-modes-holonomy-universal}) together with the classes along $e_{0}$ in the static frame (\ref{c:pairing.static-frame-only-expansion-survives}) and along $e_{0}$, $e_{1}$ in the regular one. (iii) The rotor $1 + Ae_{12}$ carries $\ket{0}$ to $\ket{0} + A\ket{12}$ (Proposition \ref{c:pairing.pairing-functional-pair-creation-amplitude}) and $\ell$ to $e_{0} + \frac{1 - A^{2}}{1 + A^{2}}e_{1} + \frac{2A}{1 + A^{2}}e_{2}$; the script verifies both states are annihilated by their rotated null pairs and computes the two directions.

\end{proof}

\begin{remark}[Departures from the registered wording]\label{c:note.pairing.departures-registered-wording}
Four, recorded in one place. (1) P2.3 was registered as ``nonzero pairing class in the $n$-$m$ planes''; computed, the class of the complex loop $(n, m)$ is zero and the nonzero one is $(n, \mb)$ (Proposition \ref{c:pairing.curvature-which-loops-move-vacuum}). The registration meant the real planes $n \wedge e_{2}$, $n \wedge e_{3}$, which do carry the class, but as worded the prediction was refuted, and it is so recorded. (2) P2.1 asks for proportionality to $\sqrt f/r$ and for a sign change across the horizon; no single frame delivers both, the first being established in the static frame (Proposition \ref{c:pairing.static-frame-only-expansion-survives}) and the second in the regular one (Proposition \ref{c:pairing.regular-frame-class-vanishes-horizon-changes}), and the two are tied by the covariance factor. (3) G1's sample point was registered as the singularity paper's necessity witness $(\tfrac35, \tfrac45)$; the computation uses its thermal sample point $e^{-\omega/T} = \tfrac{9}{16}$, $\tan\theta = \tfrac34$, for which the rotated null direction is rational; nothing depends on which. (4) The consequence drawn from K1.2 in the verdict below, that a holonomy reproducing the family mode by mode would need a loop selected by the mode, is not part of the registered condition; K1.1 and K1.2 as registered give the verdict on their own.
\end{remark}

\phantomsection\label{c:note.pairing.setting-what-does-two-claims}
\textbf{Verdict on Claim R2.} K2.1 does not trigger: the classes of Propositions \ref{c:pairing.static-frame-only-expansion-survives} and \ref{c:pairing.regular-frame-class-vanishes-horizon-changes} are nonzero off the horizon. K2.2 does not trigger: the class is $\ip{\nabla_{X}\ell}{m}$, its zeros are invariant under rescaling, and the two frames agree covariantly. K2.3 does not trigger: the classes along $\partial_{t}$ and $\partial_{r}$ vanish in both frames. K2.4 triggers: every generator found acts on the two modes of one fibre, and the object its wording asks for, a generator in $\so(V_{x} \oplus V_{y})$ pairing the completion's independent ladders at two points, is not something a connection produces. So Claim R2 is closed for the pre-ascent Dirac fibre and open for the boundary modes across points. Of the predictions, P2.2 and P2.4 held, P2.1 held with its two clauses in two frames, and P2.3 was refuted as worded (Remark \ref{c:note.pairing.departures-registered-wording}).

How much was at risk should be said plainly: nothing. K2.2 is an identity, since $\Omega$ is defined by $\nabla_{X}v = \tfrac12[\Omega(X), v]$ and $\pv$ by $\ip{B\cdot\ell}{m}$, so $\pv(\Omega(X)) = \ip{\nabla_{X}\ell}{m}$ could not have failed. K2.1 asks whether the outgoing radial congruence of Schwarzschild expands. K2.3 proposes to contradict an ascended statement about the completed pair by a pre-ascent exterior computation that cannot reach it. K2.4 was registered as expected to trigger, and did. The apparatus adjudicated nothing on this claim; what the registration did was fix the computation and its sign conventions before the numbers existed. The content of the note is the dictionary and the exact numbers.

\phantomsection\label{c:note.pairing.setting-what-does-two-claims-2}
\textbf{Verdict on Claim G1.} Not closed. K1.1 triggers: the two frames share a vacuum (Proposition \ref{c:pairing.two-frames-share-vacuum-selected-transport}(i)). K1.2 triggers: the class of the one transport the framework selects is zero (Proposition \ref{c:pairing.two-frames-share-vacuum-selected-transport}(ii)), and no single holonomy carries a mode weight. K1.3 triggers by the singularity paper's own Proposition 7.21(3): the Hawking pairing needs the Euclidean rotor, whose generator is not in the real algebra, while every holonomy is a real rotor. The registered claim ``some path or loop'' was, as worded, satisfiable by drawing a loop of any prescribed area for each mode; Proposition \ref{c:pairing.two-frames-share-vacuum-selected-transport}(iii) shows what such a family of loops would have to do, one rotor per mode, and nothing in the framework selects one, its mode weights being the s-wave weights of an imported premise \cite[Prem.~8.1]{ButlerT5} and its only selected transport the $t$-line. Two limits of this verdict are recorded. The representation of a Hawking state on the two modes of one fibre in (iii) is an analogy this note constructs, not the singularity paper's pairing, which is between a boundary mode and its partner across points. And the route by which the mode weight does enter a transport, the Killing flow continued through the Euclidean period, is exactly K1.3 and is not a holonomy of the real connection. P1.1 and P1.2 held; the regularity premise remains where the singularity paper put it.

\phantomsection\label{c:note.pairing.setting-what-does-two-claims-3}
\textbf{Prior art and the two ledgers.} Proposition \ref{c:pairing.stabiliser-parabolic-algebra-null-direction} in its geometric content is the parabolic stabiliser of a point of the celestial sphere; Proposition \ref{c:pairing.pairing-functional-pair-creation-amplitude} is the Thouless matrix element; Proposition \ref{c:pairing.pairing-class-transverse-derivative-null-direction} is the Newman-Penrose definition of the spin coefficients; Propositions \ref{c:pairing.static-frame-only-expansion-survives} and \ref{c:pairing.regular-frame-class-vanishes-horizon-changes} are the Schwarzschild spin coefficients, standard since Chandrasekhar; Proposition \ref{c:pairing.curvature-which-loops-move-vacuum} and Lemma \ref{c:pairing.loop-classes-are-weyl-components} are $\Psi_{2} = -M/r^{3}$; Proposition \ref{c:pairing.displacement-pairs-point-mode-vacuum-parallel} is twistor incidence. What the note proves for itself is the functional $\pv$ in the series' conventions, its identity with the pair-creation amplitude, the identification of the singularity paper's witness as a rotation of the null direction (Corollary \ref{c:pairing.singularity-papers-witness-rotation-null-direction}), the exact Schwarzschild classes in both frames with the horizon zero and the sign change, the curvature classes, and the three statements of Proposition \ref{c:pairing.two-frames-share-vacuum-selected-transport}. Under the two-ledger discipline every result here is transport-absent: the section ``The vacuum's stabiliser, and the pairing functional'' to the section ``The ascended reading: the point mode'' would stand verbatim with no transport principle, the pre-ascent reading using no ascent and the ascended reading reproducing incidence; what the framework supplies is the labelling of which null direction is its vacuum, and that labelling rests on the curvature pair, the sign of the square and the flag bit. The negative result is likewise transport-absent: the framework's transport does not produce the Hawking pairing. The note's contribution is a dictionary, not a derivation, and it is offered as one.

\phantomsection\label{c:note.pairing.setting-what-does-two-claims-4}
\textbf{What entanglement is, in the framework, level by level.} At one point in the pre-ascent reading, entanglement between the two modes of the fibre is the angle between two null directions, and it is produced by transport at the rate of the optical scalars of the vacuum's congruence; in Schwarzschild, for the outgoing vacuum, by the expansion alone, transversely, and not at all on the horizon. At one point in the ascended reading, it is additionally the displacement of the point against its vacuum's ray. Across points, between the ladders the completion attaches to distinct points of the exceptional set, it is produced by nothing the framework supplies; that is the pairing the Hawking pair needs, and its angle enters through the regularity premise, calibrated at a locus where the transport principle makes no statement.

\phantomsection\label{c:note.pairing.setting-what-does-two-claims-5}
\textbf{Consequences for the owning papers.} For the singularity paper, what is defensible is this. Its Premise 7.30 says the completion ``provides no generator that is off-diagonal in that basis'', the basis being the joint number basis of one ladder pair per point of $E$; nothing here touches that, since the pre-ascent connection acts on a different space and, along the $t$- and $r$-lines, lies in the stabiliser, so Theorem C.8 is untouched, and the cross-point clause is untouched because no connection produces a cross-point generator. What the note adds is a clause for Remark 13.5: the framework compares two Witt decompositions and does not derive them, and it does derive one relation between two of them, the rotation of a vacuum transported along a path, measured against the vacuum field at the endpoint, which is the path-ordered transport whose generator's pairing part is the class of Proposition \ref{c:pairing.pairing-class-transverse-derivative-null-direction}; what it does not derive is which two decompositions the Hawking pair uses, and the remark's clause on the pairing angle stands. For the laboratory paper: Proposition 4.1(3)'s null-direction datum has an exact transport law, Proposition \ref{c:pairing.pairing-class-transverse-derivative-null-direction}, and the paragraph that says the connection returns the bivector tilted can say by how much: by $\kappa, \tau, \sigma, \rho$, with the discrete bits unmoved. Successor problems: the tractor connection on the singular collar (Remark \ref{c:note.pairing.what-does-does-not-say}); a cross-point generator, which is the kinetic-term question of the singularity paper's Remark C.9 under another name; and vacuum selection, which this note does not address.

\section{Position in the series}

\phantomsection\label{c:note.pairing.setting-position-series}
This document carries no T-number: it is a companion note, deposited with T5 \cite{ButlerT5} and read with T16 \cite{ButlerT16}, and its notation is that of the primer T0 \cite{ButlerT0} and the root paper T1 \cite{ButlerT1}, from which it takes the transport principle and the forced $(1,1)$ ascent (Theorem 4.13). From T5 it takes the stabiliser-tangent split (Proposition 13.4), the two clauses of Remark 13.5, the locality premise (Premise 7.30) with its class-relative theorem (Theorem C.8) and residual (Remark C.9), the necessity witness (Corollary 11.6), the thermofield angle (Lemma 7.2, Theorem 7.27), the regularity premise and its placement (Premise 7.19, Proposition 7.25, Proposition 7.21), the $t$-line holonomy (Proposition 13.2), and the Schwarzschild connection, curvature and radial-plane selection (Appendix A, Theorem 6.1). From T16 it takes the two presentations and capacities (Propositions 3.1, 3.2), the vacuum as a null direction (Proposition 4.1(3)), Cartan's theorem (Fact 2.10) and the flat conformal model (Fact 2.3). From T3 \cite{ButlerT3} the capacity count (Theorem 4.3); from T9 \cite{ButlerT9} the incidence theorem (Theorem 4.2) and the stabiliser of $\ninf$ (Fact 4.3); from T4 \cite{ButlerT4} the ladder pair per point of the exceptional set (Theorem 8.12), which is what makes the cross-point clause a separate question; from T14 \cite{ButlerT14} the flag bit (Definition 2.10, Proposition 5.1); from T17 \cite{ButlerT17} the two-ledger grading. It sends back to T5 a sharper pairing clause for Premise 7.30 and a clause for Remark 13.5, and to T16 an exact form for Proposition 4.1(3)'s transport. The Hawking paper T7 \cite{ButlerT7} and the families paper T13 \cite{ButlerT13} share the quasi-free frontier this note does not touch; the capstone T17 \cite{ButlerT17} carries the status of the two claims.

\section{Machine verification}

\phantomsection\label{c:note.pairing.setting-machine-verification}
Every computational identity asserted in this note has been verified in exact arithmetic by the script \texttt{verify\_T05\_pairing.py} \cite{ButlerVerifyT05pairing}, which imports the series' shared library \texttt{GA.py}: its basis-blade Clifford engine over the rationals, its Gaussian rationals for the complex Witt vectors, its rational functions in one variable for the two Schwarzschild frames, and its exact ranks and kernels. Identities in $\Q(F)$ or $\Q(\beta)$ are verified as identities where the text says so, and at rational sample points where the text says so. No floating point is used anywhere; the script exits nonzero on any failure. Fifty-three checks, all passing, in eight parts. On the $(t, r, \vartheta)$ and $(T, r, \vartheta)$ blocks the check covers, the connection forms of both frames are verified against the first structure equations, which with antisymmetry determines them there, the curvature forms are derived from them by the second structure equations, and the block one-forms are tied to the frame matrices that the classes are computed from. The pre-registration is deposited beside the script \cite{ButlerPreregT05pairing}.

\emph{Part A: the stabiliser.} The tetrad's pairings; the action $\tfrac12[e_{a}e_{b}, e_{c}] = \eta_{bc}e_{a} - \eta_{ac}e_{b}$ on all basis bivectors and vectors; the six values of $\pv$ with the two transverse images $e_{02}\cdot\ell = -e_{2}$, $e_{12}\cdot\ell = e_{2}$; the equivalence of $B\cdot W \subseteq W$ with $\pv(B) = 0$ on twenty random complex bivectors, none of which stabilises, with the projection $B + \pv(B)e_{12}$ always stabilising; the null rotations $\ell e_{2}$, $\ell e_{3}$ as bivectors fixing $\ell$; the rank two of $(\operatorname{Re}\pv, \operatorname{Im}\pv)$, the rank four of $\{e_{01}, e_{23}, \ell e_{2}, \ell e_{3}\}$ and the rank five with $e_{02}$ appended; the classes $2$, $2i$ of $ne_{2}$, $ne_{3}$ and $0$, $4$ of $n m$, $n\mb$.

\emph{Part B: the Fock dictionary.} The Clifford relations of the four Jordan-Wigner operators; $c(\ell)$, $c(m)$ annihilating the vacuum and $c(n)$, $c(\mb)$ creating; $\bra{12}B\ket{0} = -\pv(B)$ on all six basis bivectors; the $2 \times 2$ structure of the even sector, including $a_{1}a_{2}\ket{12} = -\ket{0}$.

\emph{Part C: the witness.} $R\tilde R = 1$; $R\ket{0} = \tfrac35\ket{0} + \tfrac45\ket{12}$; $R\ell\tilde R$ with its nullity and its pairing $\tfrac{32}{25}$ with $\ell$; the rotated pair annihilating $R\ket{0}$; the boost rotor acting on $\ket{0}$ as $\tfrac13$.

\emph{Part D: the static frame.} $f'$ and $f''$ by differentiation in $\Q(F)$; the first structure equations on the $(t, r, \vartheta)$ block; the second structure equations giving the three curvature two-forms with their coefficients; the bivectors reproducing all sixteen frame derivatives; the block one-forms evaluated on the frame vectors equal to the frame matrices; the four classes; the scaling law with the derivative term $(e_{1}\lambda)\ell$ modelled and nonzero; the horizon zeros and the value $-\tfrac{48}{125}$; the tetrad components $(0, 0, 0, -2\sqrt f/r)$; the transverse expansion, its equality with $-\pv(\Omega(\mb))$, and $\nabla_{\ell}\ell = (f'/2\sqrt f)\ell$; $\Omega(e_{0}) \propto e_{01}$ and $\Omega(e_{1}) = 0$.

\emph{Part E: curvature.} The two eigenvalues from the derived block forms (the angular one is entered); the six classes; the null rotation $(M/r^{3})\ell e_{2}$ for the loop $(\ell, e_{2})$, the class $2M/r^{3}$ for $(n, e_{2})$, and the six complex loops, $(n, \mb)$ alone carrying $4M/r^{3}$; the Weyl dictionary: Ricci-flatness on all sixteen frame pairs, $\Psi_{0} = \Psi_{1} = 0$, $\Psi_{2} = -4M/r^{3}$, and the six loop classes equal to $(\Psi_{0}, 0, -\Psi_{1}, \Psi_{1}, 0, -\Psi_{2})$, in the adapted tetrad and again in a generic rational tetrad where $\Psi_{0}$ and $\Psi_{1}$ are nonzero.

\emph{Part F: the regular frame.} The first and second structure equations on the $(T, r, \vartheta)$ block as identities in $\Q(\beta)$, with $\beta' = -\beta^{3}/4M$ and the curvature coefficients equal to the static ones as functions of $r$; the bivectors against the frame derivatives; the block one-forms against the frame matrices; the four classes; the values $-\tfrac{16}{125}, 0, 4$ at $\beta = \tfrac45, 1, 2$; the ingoing vacuum's classes $+(1 + \beta)/r$ with the values $\tfrac{144}{125}, 2, 12$; the covariance factor $F/(1 + \beta)$ at five rational points of the circle; $\Omega(e_{0}) = 0$ and $\Omega(e_{1}) \propto e_{01}$; the transverse expansion $2(1 - \beta)/r$ with its sign on both sides and $\nabla_{\ell}\ell = (\beta^{3}/4M)\ell$; the $(n, e_{2})$ curvature class $2M/r^{3}$ in this frame, equal to the static value at three radii and to $1$ on the horizon.

\emph{Part G: gap 1.} $\Lambda^{\mathsf T}\eta\Lambda = \eta$, the images $\tfrac{F}{1 + \beta}\ell$ and $e_{2}$, and the identity on the circle at three rational points; the two Hawking states as vacua of rotated decompositions with the two distinct rational null directions.

\emph{Part H: the flat model.} $\nz$, $\ninf$; isotropy and rank of $W(x)$ at four points; the translation identity, the fixing of $\ninf$, the action on $\ell(x)$, and the second-order remainder of a finite displacement; the classes $(\ip{dx}{\ell}, \ip{dx}{m}, 0)$ at four points for twenty random real displacements; the four named real displacements and the complex $m$, $\mb$; incidence at three values of $s$; the rank three of the three functionals on the fifteen bivectors, their common kernel of dimension twelve, and its coincidence with the stabiliser on ten random kernel elements and on every basis bivector.

What the script does not check: the geometric identifications imported as Facts \ref{c:pairing.imported-null-flags}, \ref{c:pairing.imported-spin-coefficients-null-tetrad}, \ref{c:pairing.imported-flat-conformal-model-twistor-incidence} and \ref{c:pairing.imported-weyl-scalars-petrov-type-schwarzschild} (the note computes components and matches them to named quantities; it does not re-derive the formalisms), Lemma \ref{c:pairing.loop-classes-are-weyl-components} beyond its Schwarzschild instance; the angular-plane curvature coefficient $(1 - f)/r^{2}$, which lies off the $(t, r, \vartheta)$ block and is taken from the singularity paper's verified appendix; and the $\hat3$ forms of both frames beyond the mirror symmetry, the $e_{23}$ term being in the stabiliser by Part A.

\section{Results used from elsewhere}
The proofs above use the following results of the series, stated briefly; each is proved in the paper named.

\begin{fact}[Imported: stabilisers in the flat model]\label{c:correspondence.imported-stabilisers-flat-model}
In $\mathfrak{so}(2,4)$ the stabiliser of the null vector at infinity is the Poincar\'e algebra, the stabiliser of its ray is the eleven-dimensional maximal parabolic, and its nilradical is the translations, which are the no-go theorem's transvections.
\emph{Source:} T9, Fact 4.3 (\href{https://doi.org/10.5281/zenodo.22722290}{doi:10.5281/zenodo.22722290}).
\end{fact}

\begin{fact}[The incidence relation is the chart]\label{c:correspondence.incidence-relation-chart}
The Pl\"ucker bivector of the incidence line of a map $A$ is exactly the completion's null point embedding $X(m(A))$, so Penrose's incidence relation is the completion's affine chart, with a determinant as its quadratic term.
\emph{Source:} T9, Thm 4.2 (\href{https://doi.org/10.5281/zenodo.22722290}{doi:10.5281/zenodo.22722290}).
\end{fact}

\begin{fact}[Imported: Cartan's purity theorem and the Thouless form]\label{c:one-relation.imported-cartans-purity-theorem-thouless-form}
Imported from Cartan, Chevalley and Thouless: pure spinors are Witt vacua forming one orbit of dimension $\nu(\nu-1)/2$, all Weyl spinors are pure for $2\nu \leq 6$, and vacua overlapping $\ket{0}$ have the exponential Thouless form.
\emph{Source:} 'E. Cartan, Lecons sur la th'eorie des spineurs, Hermann (1938).; C. Chevalley, The Algebraic Theory of Spinors, Columbia University Press (1954).; D. J. Thouless, Stability conditions and nuclear rotations in the Hartree-Fock theory, Nucl. Phys. 21 (1960) 225--232..
\end{fact}

\begin{fact}[Imported: pure spinors and the incidence chart]\label{c:one-relation.imported-pure-spinors-incidence-chart}
Imported: every conformal spinor is a pure spinor, the Fock vacuum of its own three-pair system, and spacetime points are null vectors of $\C^6$, equivalently lines of the projective spinor space.
\emph{Source:} T16, Fact 2.3 (\href{https://doi.org/10.5281/zenodo.22723212}{doi:10.5281/zenodo.22723212}).
\end{fact}

\begin{fact}[Spacetime as the manifold of vacuum lines]\label{c:one-relation.spacetime-manifold-vacuum-lines}
Every spacetime point is a null vector and so a projective line of Fock vacua, and two points are null-separated exactly when their lines share a vacuum.
\emph{Source:} T16, Prop 4.1 (\href{https://doi.org/10.5281/zenodo.22723212}{doi:10.5281/zenodo.22723212}).
\end{fact}

\begin{fact}[The two capacities, and three consequences]\label{c:one-relation.two-capacities-three-consequences}
A real quadratic space of signature $(p,q)$ has real capacity $\min(p,q)$ and complex capacity $\lfloor (p+q)/2 \rfloor$, both rising by one per ascent, giving one boost qubit for $(1,3)$ and complex capacity three for $(2,4)$.
\emph{Source:} T16, Prop 3.2 (\href{https://doi.org/10.5281/zenodo.22723212}{doi:10.5281/zenodo.22723212}).
\end{fact}

\begin{fact}[The two presentations]\label{c:one-relation.two-presentations}
Read on an orthogonal basis the Clifford relation is a metric, and read on a Witt basis it is the canonical anticommutation relations, so the geometric and quantum readings are two charts on one algebra.
\emph{Source:} T16, Prop 3.1 (\href{https://doi.org/10.5281/zenodo.22723212}{doi:10.5281/zenodo.22723212}).
\end{fact}

\begin{fact}[The curvature of the gauge field: exact shape decoupling]\label{c:singularity.curvature-gauge-field-exact-shape-decoupling}
The curvature of the black-hole gauge field is exactly $\mathcal{R} = -(M/r^3)\,S$ with $S = 1 - 3\Pi$ a fixed shape operator of spectrum $(-2,-2,1,1,1,1)$.
\emph{Source:} T5, Thm 6.1 (\href{https://doi.org/10.5281/zenodo.22721790}{doi:10.5281/zenodo.22721790}).
\end{fact}

\begin{fact}[The envelope and the endpoint, and why boundary-mode locality is needed for them]\label{c:singularity.envelope-endpoint-why-boundary-mode-locality}
Under (L), (U), (C), (I) and (L$_E$), $S_R(t) \leq \min(\Sth(t), C(t))$ with $S_R = 0$ at complete evaporation, and the bound can fail without (L$_E$).
\emph{Source:} T5, Cor 11.6 (\href{https://doi.org/10.5281/zenodo.22721790}{doi:10.5281/zenodo.22721790}).
\end{fact}

\begin{fact}[The Euclidean rotor, and where it lies]\label{c:singularity.euclidean-rotor-where-it-lies}
The imaginary-rapidity continuation $R_{i\theta} = \cos(\theta/2) + i\sigma\sin(\theta/2)$ of the radial boost rotor has period $4\pi$, with conjugation period $2\pi$, and is realised by no real rotor of $\Cl(1,3)$, in particular not by the curvature-selected pair, whose dual member $I\sigma_r = e_2e_3$ rotates along the bifurcation sphere.
\emph{Source:} T5, Prop 7.21 (\href{https://doi.org/10.5281/zenodo.22721790}{doi:10.5281/zenodo.22721790}).
\end{fact}

\begin{fact}[(L$_E$) as the stabiliser and the tangent]\label{c:singularity.l-stabiliser-tangent}
In $\mathfrak{so}(2\nu)$ on Fock space, hops and pair annihilations stabilise the vacuum ray and pair creations span the tangent to its orbit, so only pairing can move the vacuum.
\emph{Source:} T5, Prop 13.4 (\href{https://doi.org/10.5281/zenodo.22721790}{doi:10.5281/zenodo.22721790}).
\end{fact}

\begin{fact}[The modular flow of the boundary modes as a holonomy, and its universal temperature]\label{c:singularity.modular-flow-boundary-modes-holonomy-universal}
The $t$-line transport in the interior collar is a boost in the completed pair's plane whose rate is $\kappa$ at the horizon, matching the modular flow at the period of (R), and whose Hubble-normalised limit at the exceptional set gives $\beta_{\mathrm{H}} = 6\pi$ if the KMS state extends to the collar, and the non-commuting limiting boosts leave $B_1$, the $t$-direction, as the one the modular flow selects.
\emph{Source:} T5, Prop 13.2 (\href{https://doi.org/10.5281/zenodo.22721790}{doi:10.5281/zenodo.22721790}).
\end{fact}

\begin{fact}[Bivectors form the Lie algebra of the orthogonal group]\label{c:std-algebra.bivector-lie-algebra}
$\Cl^{2}$ is closed under $A \times C = \langle AC \rangle_2$, acts skew-adjointly on $V$ by $A \times x = \langle Ax \rangle_1$, and is isomorphic to $\so(V,Q)$ when $Q$ is nondegenerate.
\emph{Source:} the algebra core of the Transport Lab, proved there with an exact check; standard, see C. Doran and A. Lasenby, Geometric Algebra for Physicists, Cambridge University Press (2003), Ch. 4 (foundations of geometric algebra: rotors, versors).
\end{fact}

\begin{fact}[Contraction by a vector is a derivation]\label{c:std-algebra.contraction-expansion}
$a \rfloor$ is a graded derivation of both products with $(a \rfloor)^{2} = 0$, so $a \rfloor (b_1 \wedge \dots \wedge b_k) = \sum_i (-1)^{i-1} B(a, b_i)\, b_1 \wedge \dots \check{b}_i \dots \wedge b_k$.
\emph{Source:} the algebra core of the Transport Lab, proved there with an exact check; standard, see D. Hestenes and G. Sobczyk, Clifford Algebra to Geometric Calculus, Reidel, Dordrecht (1984), Ch. 1 (geometric algebra).
\end{fact}

\begin{fact}[The Fock representation, and the Fock space as a minimal left ideal]\label{c:std-algebra.fock-representation}
$\Cl(\R^{k,k}) \cong \End(\Lambda W)$ by wedge and contraction, and the Fock space is the minimal left ideal $\Cl\, f$ of the vacuum idempotent $f = \prod_i v_i w_i$.
\emph{Source:} the algebra core of the Transport Lab, proved there with an exact check; standard, see C. Chevalley, The Algebraic Theory of Spinors, Columbia University Press (1954), Ch. II and III (spinors, pure spinors).
\end{fact}

\begin{fact}[Orthogonal vectors anticommute]\label{c:std-algebra.orthogonal-anticommute}
$uv = -vu$ iff $B(u,v) = 0$; in an orthogonal basis $e_I e_J = \pm \prod_{I \cap J} Q(e_i)\; e_{I \triangle J}$ with the sorting sign.
\emph{Source:} the algebra core of the Transport Lab, proved there with an exact check; standard, see C. Doran and A. Lasenby, Geometric Algebra for Physicists, Cambridge University Press (2003), Ch. 4 (foundations of geometric algebra).
\end{fact}

\begin{fact}[Imported: the dictionary between Gauge Theory Gravity and curved-spacetime language]\label{c:synthesis.imported-gtg-dictionary}
On the invertible region, tangent vectors are $\hb^{-1}(a)$, the Levi-Civita derivative is $D_a$, geodesics have $D$-parallel velocity, and the Riemann, Ricci, Schouten and Weyl tensors and Killing fields are read from the field strength $\mathcal{R}(B)$ and $D$.
\emph{Source:} A. Lasenby, C. Doran and S. Gull, Gravity, gauge theories and geometric algebra, Phil. Trans. R. Soc. Lond. A 356 (1998), 487--582., \S\S 3.5, 4.3--4.4, 5.1--5.2, App.~C; C. Doran and A. Lasenby, Geometric Algebra for Physicists, Cambridge University Press (2003)., \S\S~13.5--13.6, 14.1; A. N. Lasenby, Geometric algebra, gravity and gravitational waves, Adv. Appl. Clifford Algebras (2019)., \S\S~4--5, eq.~(8.5) (arXiv:1912.05960); R. M. Wald, General Relativity, University of Chicago Press (1984)., \S~12.5.
\end{fact}

\begin{fact}[Witt index is quantum capacity]\label{c:tower.witt-index-quantum-capacity}
For a nondegenerate real quadratic space of signature $(p,q)$, the maximal number of fermionic mode pairs realisable by vectors of its Clifford algebra is the Witt index $\min(p,q)$, which rises by exactly one at each rung of the tower.
\emph{Source:} T3, Thm 4.3 (\href{https://doi.org/10.5281/zenodo.22721556}{doi:10.5281/zenodo.22721556}).
\end{fact}

\begin{thebibliography}{99}
\bibitem{ButlerT5} L.~Butler,
\emph{Division by Zero at the Gravitational Singularity: The Transport Principle Applied to the Black Hole, and the Information Paradox as a Category Error},
Transport Papers T5 (2026), doi:\href{https://doi.org/10.5281/zenodo.22721790}{10.5281/zenodo.22721790}.
\bibitem{ButlerT16} L.~Butler,
\emph{Quantum Mechanics and General Relativity from One Relation: The Two Shadows and the Laboratory},
Transport Papers T16 (2026), doi:\href{https://doi.org/10.5281/zenodo.22723212}{10.5281/zenodo.22723212}.
\bibitem{ButlerT1} L.~Butler,
\emph{The Transport Principle: Total Division in Clifford Algebras via Minimal Completion},
Transport Papers T1 (2026).
\bibitem{ButlerT4} L.~Butler,
\emph{Clifford Sheaves over Degenerate Metrics: Nilpotent Ideals, One-Sided Completion, and the Exceptional Algebra},
Transport Papers T4 (2026), doi:\href{https://doi.org/10.5281/zenodo.22721669}{10.5281/zenodo.22721669}.
\bibitem{ButlerT14} L.~Butler,
\emph{Dark Matter and the Arrow of Time: The Five Dark Towers},
Transport Papers T14 (2026).
\bibitem{ButlerT17} L.~Butler,
\emph{Capstone: The Transport Programme, Method, Status, and the Falsifiable Ledger},
Transport Papers T17 (2026).
\bibitem{Cartan1938} \'E.~Cartan,
\emph{Le\c{c}ons sur la th\'eorie des spineurs},
Hermann, Paris (1938).
\bibitem{Chevalley1954} C.~Chevalley,
\emph{The Algebraic Theory of Spinors},
Columbia University Press (1954).
\bibitem{Thouless1960} D.~J. Thouless,
\emph{Stability conditions and nuclear rotations in the Hartree-Fock theory},
Nucl. Phys. 21 (1960) 225--232.
\bibitem{ButlerT3} L.~Butler,
\emph{The Tower and the CAR Limit: Iterated Transport, the Frozen Clock, and Second Quantisation},
Transport Papers T3 (2026), doi:\href{https://doi.org/10.5281/zenodo.22721556}{10.5281/zenodo.22721556}.
\bibitem{PenroseRindler1} R.~Penrose, W.~Rindler,
\emph{Spinors and Space-Time, Vol.~1: Two-Spinor Calculus and Relativistic Fields},
Cambridge University Press (1984).
\bibitem{NewmanPenrose1962} E.~Newman, R.~Penrose,
\emph{An approach to gravitational radiation by a method of spin coefficients},
J. Math. Phys. 3 (1962) 566--578.
\bibitem{Wald1984} R.~M. Wald,
\emph{General Relativity},
University of Chicago Press (1984).
\bibitem{ButlerT9} L.~Butler,
\emph{The Twistor Correspondence: Ascended Spinors, Incidence, and the Palatial Bridge},
Transport Papers T9 (2026).
\bibitem{BEG1994} T.~N.~Bailey, M.~G.~Eastwood and A.~R.~Gover,
Thomas's structure bundle for conformal, projective and related structures,
\emph{Rocky Mountain J.\ Math.} \textbf{24} (1994), 1191--1217.
\bibitem{ButlerPreregT05pairing} L.~Butler,
\emph{preregistration\_T05\_pairing\_type.md: the kill conditions and predictions of the pairing-creator experiment of ``The Descent, Lifted'', as registered before computation},
deposited with this note (2026).
\bibitem{LDG1998} A.~Lasenby, C.~Doran and S.~Gull,
Gravity, gauge theories and geometric algebra,
\emph{Phil.\ Trans.\ R.\ Soc.\ Lond.\ A} \textbf{356} (1998), 487--582.
\bibitem{PenroseRindler2} R.~Penrose, W.~Rindler,
\emph{Spinors and Space-Time, Vol.~2: Spinor and Twistor Methods in Space-Time Geometry},
Cambridge University Press (1986).
\bibitem{ButlerT0} L.~Butler,
\emph{An Introduction to Geometric Algebra: A Primer for the Transport Papers},
Transport Papers T0 (2026).
\bibitem{ButlerT7} L.~Butler,
\emph{Hawking Radiation from the Ledger: State, Continuum, Lift, and the Finite-Ledger Seam},
Transport Papers T7 (2026), doi:\href{https://doi.org/10.5281/zenodo.22722121}{10.5281/zenodo.22722121}.
\bibitem{ButlerT13} L.~Butler,
\emph{Quark and Lepton Families: Condensation, Democracy, and the Koide Relation},
Transport Papers T13 (2026), doi:\href{https://doi.org/10.5281/zenodo.22722845}{10.5281/zenodo.22722845}.
\bibitem{ButlerVerifyT05pairing} L.~Butler,
\emph{verify\_T05\_pairing.py: exact-arithmetic verification of the computational identities of ``Where the Pairing Generator Is''},
Python script, deposited with this note alongside the shared library GA.py (2026).
\end{thebibliography}
\end{document}